%%%%%%%%%%%%%%%%%%%%%%%%%%%%%%%%%%%%%%%%%%%%%%%%%%%%%%%%%%%%
% 11.2 NUMERICAL LINEAR ALGEBRA
% The Composition of Advanced Mathematics
%%%%%%%%%%%%%%%%%%%%%%%%%%%%%%%%%%%%%%%%%%%%%%%%%%%%%%%%%%%%

\section{Numerical Linear Algebra}
\label{sec:numerical-linear-algebra}

\prerequisite{
Linear algebra, matrices, determinants, vector spaces, eigenvalues,
orthogonality, least squares, numerical analysis, floating-point
arithmetic, conditioning, and basic algorithmic complexity.
}

\learningobjectives{
By the end of this section, the reader should be able to:
\begin{itemize}
    \item explain the goals of numerical linear algebra;
    \item distinguish dense and sparse matrix problems;
    \item interpret matrix and vector norms;
    \item understand induced matrix norms;
    \item define the condition number of a linear system;
    \item distinguish residual, forward error, and backward error;
    \item understand why Gaussian elimination requires pivoting;
    \item perform LU factorization;
    \item understand partial and complete pivoting;
    \item formulate triangular solves;
    \item understand computational complexity of direct linear solvers;
    \item recognize Cholesky factorization;
    \item identify symmetric positive-definite matrices;
    \item understand QR factorization;
    \item apply Gram--Schmidt conceptually;
    \item distinguish classical and modified Gram--Schmidt;
    \item understand Householder reflections;
    \item understand Givens rotations;
    \item solve least-squares problems using QR factorization;
    \item explain why normal equations can worsen conditioning;
    \item understand singular value decomposition;
    \item define singular values;
    \item understand matrix rank and numerical rank;
    \item compute pseudoinverse solutions conceptually;
    \item understand low-rank approximation;
    \item state the Eckart--Young theorem conceptually;
    \item understand eigenvalue computation;
    \item apply the power method;
    \item understand inverse iteration;
    \item understand Rayleigh quotient iteration conceptually;
    \item understand QR iteration conceptually;
    \item recognize Hessenberg and tridiagonal reductions;
    \item distinguish direct and iterative linear solvers;
    \item understand Jacobi iteration;
    \item understand Gauss--Seidel iteration;
    \item understand successive over-relaxation conceptually;
    \item formulate stationary iterations using matrix splittings;
    \item understand spectral-radius convergence conditions;
    \item recognize Krylov subspaces;
    \item understand conjugate-gradient methods;
    \item recognize GMRES conceptually;
    \item understand preconditioning;
    \item recognize sparse matrix storage;
    \item understand fill-in;
    \item understand iterative refinement;
    \item distinguish algebraic error from floating-point error;
    \item understand backward stability of linear-system solvers;
    \item connect numerical linear algebra with optimization,
          differential equations, statistics, data science, and
          scientific computing.
\end{itemize}
}

%%%%%%%%%%%%%%%%%%%%%%%%%%%%%%%%%%%%%%%%%%%%%%%%%%%%%%%%%%%%
\subsection{What Is Numerical Linear Algebra?}
\label{subsec:what-is-numerical-linear-algebra}
%%%%%%%%%%%%%%%%%%%%%%%%%%%%%%%%%%%%%%%%%%%%%%%%%%%%%%%%%%%%

Numerical linear algebra develops reliable and efficient algorithms for
matrix computations.

Core problems include:

\[
\boxed{
A\mathbf{x}
=
\mathbf{b},
}
\]

\[
\boxed{
A\mathbf{x}
\approx
\mathbf{b},
}
\]

\[
\boxed{
A\mathbf{v}
=
\lambda\mathbf{v},
}
\]

and

\[
\boxed{
A
\approx
\text{low-rank matrix}.
}
\]

These problems appear throughout applied mathematics.

%%%%%%%%%%%%%%%%%%%%%%%%%%%%%%%%%%%%%%%%%%%%%%%%%%%%%%%%%%%%
\subsection{Why Numerical Linear Algebra Matters}
\label{subsec:why-numerical-linear-algebra-matters}
%%%%%%%%%%%%%%%%%%%%%%%%%%%%%%%%%%%%%%%%%%%%%%%%%%%%%%%%%%%%

Many computational models eventually require one or more matrix
operations.

Examples include:

\[
\boxed{
\text{Newton steps in optimization},
}
\]

\[
\boxed{
\text{finite-difference PDE systems},
}
\]

\[
\boxed{
\text{least-squares regression},
}
\]

\[
\boxed{
\text{Markov-chain calculations},
}
\]

and

\[
\boxed{
\text{data compression}.
}
\]

The efficiency and stability of the matrix computation can dominate the
overall computation.

%%%%%%%%%%%%%%%%%%%%%%%%%%%%%%%%%%%%%%%%%%%%%%%%%%%%%%%%%%%%
\subsection{Dense and Sparse Matrices}
\label{subsec:dense-sparse-matrices}
%%%%%%%%%%%%%%%%%%%%%%%%%%%%%%%%%%%%%%%%%%%%%%%%%%%%%%%%%%%%

A dense matrix has many nonzero entries.

A sparse matrix has relatively few nonzero entries.

For a matrix

\[
A\in\mathbb{R}^{m\times n},
\]

a dense representation stores approximately

\[
\boxed{
mn
}
\]

entries.

For sparse matrices, storing every zero entry may be wasteful.

%%%%%%%%%%%%%%%%%%%%%%%%%%%%%%%%%%%%%%%%%%%%%%%%%%%%%%%%%%%%
\subsection{Vector Norms}
\label{subsec:numerical-vector-norms}
%%%%%%%%%%%%%%%%%%%%%%%%%%%%%%%%%%%%%%%%%%%%%%%%%%%%%%%%%%%%

Important vector norms include:

\[
\boxed{
\|\mathbf{x}\|_1
=
\sum_{i=1}^{n}
|x_i|,
}
\]

\[
\boxed{
\|\mathbf{x}\|_2
=
\sqrt{
\sum_{i=1}^{n}
x_i^2
},
}
\]

and

\[
\boxed{
\|\mathbf{x}\|_\infty
=
\max_i
|x_i|.
}
\]

These norms measure vector size in different ways.

%%%%%%%%%%%%%%%%%%%%%%%%%%%%%%%%%%%%%%%%%%%%%%%%%%%%%%%%%%%%
\subsection{Induced Matrix Norms}
\label{subsec:induced-matrix-norms}
%%%%%%%%%%%%%%%%%%%%%%%%%%%%%%%%%%%%%%%%%%%%%%%%%%%%%%%%%%%%

A vector norm induces a matrix norm defined by

\[
\boxed{
\|A\|
=
\max_{\mathbf{x}\neq0}
\frac{
\|A\mathbf{x}\|
}{
\|\mathbf{x}\|
}.
}
\]

Important examples include:

\[
\boxed{
\|A\|_1
=
\max_j
\sum_i
|a_{ij}|,
}
\]

and

\[
\boxed{
\|A\|_\infty
=
\max_i
\sum_j
|a_{ij}|.
}
\]

Thus the \(1\)-norm is the maximum absolute column sum, while the
\(\infty\)-norm is the maximum absolute row sum.

%%%%%%%%%%%%%%%%%%%%%%%%%%%%%%%%%%%%%%%%%%%%%%%%%%%%%%%%%%%%
\subsection{Matrix \(2\)-Norm}
\label{subsec:matrix-two-norm}
%%%%%%%%%%%%%%%%%%%%%%%%%%%%%%%%%%%%%%%%%%%%%%%%%%%%%%%%%%%%

The matrix \(2\)-norm satisfies

\[
\boxed{
\|A\|_2
=
\sigma_{\max}(A),
}
\]

where

\[
\sigma_{\max}(A)
\]

is the largest singular value of \(A\).

This norm plays an important role in conditioning and least squares.

%%%%%%%%%%%%%%%%%%%%%%%%%%%%%%%%%%%%%%%%%%%%%%%%%%%%%%%%%%%%
\subsection{Frobenius Norm}
\label{subsec:frobenius-norm}
%%%%%%%%%%%%%%%%%%%%%%%%%%%%%%%%%%%%%%%%%%%%%%%%%%%%%%%%%%%%

The Frobenius norm is

\[
\boxed{
\|A\|_F
=
\sqrt{
\sum_{i=1}^{m}
\sum_{j=1}^{n}
|a_{ij}|^2
}.
}
\]

It can also be written

\[
\boxed{
\|A\|_F
=
\sqrt{
\operatorname{tr}(A^TA)
}.
}
\]

Unlike induced norms, it treats the matrix entries collectively as a
vector.

%%%%%%%%%%%%%%%%%%%%%%%%%%%%%%%%%%%%%%%%%%%%%%%%%%%%%%%%%%%%
\subsection{Linear Systems}
\label{subsec:numerical-linear-systems}
%%%%%%%%%%%%%%%%%%%%%%%%%%%%%%%%%%%%%%%%%%%%%%%%%%%%%%%%%%%%

Consider

\[
\boxed{
A\mathbf{x}
=
\mathbf{b},
}
\]

where

\[
A\in\mathbb{R}^{n\times n}.
\]

If \(A\) is nonsingular, the exact solution is formally

\[
\mathbf{x}
=
A^{-1}\mathbf{b}.
\]

Numerically, one generally does not compute \(A^{-1}\) explicitly.

Instead, one solves the system using factorization or iterative methods.

%%%%%%%%%%%%%%%%%%%%%%%%%%%%%%%%%%%%%%%%%%%%%%%%%%%%%%%%%%%%
\subsection{Why Explicit Matrix Inversion Is Usually Avoided}
\label{subsec:avoid-explicit-inverse}
%%%%%%%%%%%%%%%%%%%%%%%%%%%%%%%%%%%%%%%%%%%%%%%%%%%%%%%%%%%%

Computing

\[
A^{-1}
\]

explicitly is typically unnecessary when the goal is only to solve

\[
A\mathbf{x}
=
\mathbf{b}.
\]

Factorization-based solves are usually:

\[
\boxed{
\text{more efficient},
}
\]

\[
\boxed{
\text{more stable},
}
\]

and

\[
\boxed{
\text{easier to reuse for multiple right-hand sides}.
}
\]

%%%%%%%%%%%%%%%%%%%%%%%%%%%%%%%%%%%%%%%%%%%%%%%%%%%%%%%%%%%%
\subsection{Residual of a Linear System}
\label{subsec:linear-system-residual}
%%%%%%%%%%%%%%%%%%%%%%%%%%%%%%%%%%%%%%%%%%%%%%%%%%%%%%%%%%%%

For an approximate solution

\[
\widehat{\mathbf{x}},
\]

the residual is

\[
\boxed{
\mathbf{r}
=
\mathbf{b}
-
A\widehat{\mathbf{x}}.
}
\]

If

\[
\mathbf{r}=0,
\]

then the approximation exactly satisfies the original linear equations.

A small residual alone does not guarantee a small forward error.

%%%%%%%%%%%%%%%%%%%%%%%%%%%%%%%%%%%%%%%%%%%%%%%%%%%%%%%%%%%%
\subsection{Forward Error in Linear Systems}
\label{subsec:linear-system-forward-error}
%%%%%%%%%%%%%%%%%%%%%%%%%%%%%%%%%%%%%%%%%%%%%%%%%%%%%%%%%%%%

If \(\mathbf{x}\) is the exact solution, the forward error is

\[
\boxed{
\mathbf{e}
=
\mathbf{x}
-
\widehat{\mathbf{x}}.
}
\]

Because

\[
A\mathbf{x}
=
\mathbf{b},
\]

and

\[
A\widehat{\mathbf{x}}
=
\mathbf{b}
-
\mathbf{r},
\]

we have

\[
\boxed{
A\mathbf{e}
=
\mathbf{r}.
}
\]

Thus,

\[
\boxed{
\mathbf{e}
=
A^{-1}\mathbf{r}.
}
\]

%%%%%%%%%%%%%%%%%%%%%%%%%%%%%%%%%%%%%%%%%%%%%%%%%%%%%%%%%%%%
\subsection{Condition Number}
\label{subsec:matrix-condition-number}
%%%%%%%%%%%%%%%%%%%%%%%%%%%%%%%%%%%%%%%%%%%%%%%%%%%%%%%%%%%%

For nonsingular \(A\), the condition number in a chosen induced norm is

\[
\boxed{
\kappa(A)
=
\|A\|
\|A^{-1}\|.
}
\]

Since

\[
\|I\|
\leq
\|A\|
\|A^{-1}\|,
\]

we have

\[
\boxed{
\kappa(A)\geq1.
}
\]

Large condition numbers indicate sensitivity.

%%%%%%%%%%%%%%%%%%%%%%%%%%%%%%%%%%%%%%%%%%%%%%%%%%%%%%%%%%%%
\subsection{Condition Number in the \(2\)-Norm}
\label{subsec:condition-number-two-norm}
%%%%%%%%%%%%%%%%%%%%%%%%%%%%%%%%%%%%%%%%%%%%%%%%%%%%%%%%%%%%

For nonsingular \(A\),

\[
\boxed{
\kappa_2(A)
=
\frac{
\sigma_{\max}(A)
}{
\sigma_{\min}(A)
}.
}
\]

If the smallest singular value is close to zero, then

\[
\kappa_2(A)
\]

is large.

Thus nearly singular matrices are ill conditioned.

%%%%%%%%%%%%%%%%%%%%%%%%%%%%%%%%%%%%%%%%%%%%%%%%%%%%%%%%%%%%
\subsection{Relative Error Bound}
\label{subsec:linear-relative-error-bound}
%%%%%%%%%%%%%%%%%%%%%%%%%%%%%%%%%%%%%%%%%%%%%%%%%%%%%%%%%%%%

For perturbations in the right-hand side only, a standard estimate is

\[
\boxed{
\frac{
\|\Delta\mathbf{x}\|
}{
\|\mathbf{x}\|
}
\leq
\kappa(A)
\frac{
\|\Delta\mathbf{b}\|
}{
\|\mathbf{b}\|
}
}
\]

under compatible assumptions.

Thus the condition number bounds the amplification of relative input
errors.

%%%%%%%%%%%%%%%%%%%%%%%%%%%%%%%%%%%%%%%%%%%%%%%%%%%%%%%%%%%%
\subsection{Worked Example: Ill Conditioning}
\label{subsec:worked-ill-conditioned-system}
%%%%%%%%%%%%%%%%%%%%%%%%%%%%%%%%%%%%%%%%%%%%%%%%%%%%%%%%%%%%

\begin{workedexamplebox}{Nearly Dependent Equations}

Consider

\[
\begin{aligned}
x+y&=2,\\
x+1.0001y&=2.0001.
\end{aligned}
\]

The equations are nearly parallel.

Small perturbations in the right-hand sides can produce relatively large
changes in the intersection point.

The matrix

\[
A
=
\begin{pmatrix}
1&1\\
1&1.0001
\end{pmatrix}
\]

is close to singular because its rows are nearly linearly dependent.

\begin{answerbox}
\[
\boxed{
\text{The system is sensitive because }A\text{ is nearly singular.}
}
\]
\end{answerbox}

\end{workedexamplebox}

%%%%%%%%%%%%%%%%%%%%%%%%%%%%%%%%%%%%%%%%%%%%%%%%%%%%%%%%%%%%
\subsection{Gaussian Elimination}
\label{subsec:gaussian-elimination-numerical}
%%%%%%%%%%%%%%%%%%%%%%%%%%%%%%%%%%%%%%%%%%%%%%%%%%%%%%%%%%%%

Gaussian elimination transforms

\[
A\mathbf{x}
=
\mathbf{b}
\]

into an equivalent upper-triangular system.

At step \(k\), multiples of row \(k\) are subtracted from rows below it to
eliminate entries

\[
a_{ik},
\qquad
i>k.
\]

The multipliers are

\[
\boxed{
m_{ik}
=
\frac{
a_{ik}
}{
a_{kk}
}.
}
\]

%%%%%%%%%%%%%%%%%%%%%%%%%%%%%%%%%%%%%%%%%%%%%%%%%%%%%%%%%%%%
\subsection{Back Substitution}
\label{subsec:back-substitution}
%%%%%%%%%%%%%%%%%%%%%%%%%%%%%%%%%%%%%%%%%%%%%%%%%%%%%%%%%%%%

After elimination, solve

\[
U\mathbf{x}
=
\mathbf{c},
\]

where \(U\) is upper triangular.

Starting with the final equation,

\[
\boxed{
x_n
=
\frac{
c_n
}{
u_{nn}
}.
}
\]

Then for

\[
i=n-1,\ldots,1,
\]

\[
\boxed{
x_i
=
\frac{
c_i
-
\sum_{j=i+1}^{n}
u_{ij}x_j
}{
u_{ii}
}.
}
\]

%%%%%%%%%%%%%%%%%%%%%%%%%%%%%%%%%%%%%%%%%%%%%%%%%%%%%%%%%%%%
\subsection{Forward Substitution}
\label{subsec:forward-substitution}
%%%%%%%%%%%%%%%%%%%%%%%%%%%%%%%%%%%%%%%%%%%%%%%%%%%%%%%%%%%%

For a lower-triangular system

\[
L\mathbf{y}
=
\mathbf{b},
\]

solve from the first equation downward.

If \(L\) has nonzero diagonal entries,

\[
\boxed{
y_i
=
\frac{
b_i
-
\sum_{j=1}^{i-1}
\ell_{ij}y_j
}{
\ell_{ii}
}.
}
\]

%%%%%%%%%%%%%%%%%%%%%%%%%%%%%%%%%%%%%%%%%%%%%%%%%%%%%%%%%%%%
\subsection{LU Factorization}
\label{subsec:lu-factorization}
%%%%%%%%%%%%%%%%%%%%%%%%%%%%%%%%%%%%%%%%%%%%%%%%%%%%%%%%%%%%

Gaussian elimination can be expressed as

\[
\boxed{
A
=
LU,
}
\]

where:

\[
L
=
\text{lower triangular matrix},
\]

and

\[
U
=
\text{upper triangular matrix}.
\]

Then

\[
A\mathbf{x}
=
\mathbf{b}
\]

becomes

\[
LU\mathbf{x}
=
\mathbf{b}.
\]

First solve

\[
\boxed{
L\mathbf{y}
=
\mathbf{b},
}
\]

then solve

\[
\boxed{
U\mathbf{x}
=
\mathbf{y}.
}
\]

%%%%%%%%%%%%%%%%%%%%%%%%%%%%%%%%%%%%%%%%%%%%%%%%%%%%%%%%%%%%
\subsection{Worked Example: LU Factorization}
\label{subsec:worked-lu-factorization}
%%%%%%%%%%%%%%%%%%%%%%%%%%%%%%%%%%%%%%%%%%%%%%%%%%%%%%%%%%%%

\begin{workedexamplebox}{Factoring a \(2\times2\) Matrix}

Let

\[
A
=
\begin{pmatrix}
2&1\\
6&5
\end{pmatrix}.
\]

The elimination multiplier is

\[
m_{21}
=
\frac62
=
3.
\]

Subtract \(3\) times row \(1\) from row \(2\):

\[
U
=
\begin{pmatrix}
2&1\\
0&2
\end{pmatrix}.
\]

The multiplier is stored in

\[
L
=
\begin{pmatrix}
1&0\\
3&1
\end{pmatrix}.
\]

Therefore,

\begin{answerbox}
\[
\boxed{
A
=
LU
=
\begin{pmatrix}
1&0\\
3&1
\end{pmatrix}
\begin{pmatrix}
2&1\\
0&2
\end{pmatrix}.
}
\]
\end{answerbox}

\end{workedexamplebox}

%%%%%%%%%%%%%%%%%%%%%%%%%%%%%%%%%%%%%%%%%%%%%%%%%%%%%%%%%%%%
\subsection{Pivoting}
\label{subsec:numerical-pivoting}
%%%%%%%%%%%%%%%%%%%%%%%%%%%%%%%%%%%%%%%%%%%%%%%%%%%%%%%%%%%%

Gaussian elimination can become unstable if a pivot

\[
a_{kk}
\]

is zero or very small.

Pivoting changes row or column order to select more appropriate pivots.

The most common strategy is partial pivoting.

%%%%%%%%%%%%%%%%%%%%%%%%%%%%%%%%%%%%%%%%%%%%%%%%%%%%%%%%%%%%
\subsection{Partial Pivoting}
\label{subsec:partial-pivoting}
%%%%%%%%%%%%%%%%%%%%%%%%%%%%%%%%%%%%%%%%%%%%%%%%%%%%%%%%%%%%

At elimination step \(k\), partial pivoting selects the entry of largest
absolute value in column \(k\) among rows \(k,\ldots,n\).

Thus choose row \(p\) satisfying

\[
\boxed{
|a_{pk}|
=
\max_{i\geq k}
|a_{ik}|.
}
\]

Then interchange rows \(k\) and \(p\).

%%%%%%%%%%%%%%%%%%%%%%%%%%%%%%%%%%%%%%%%%%%%%%%%%%%%%%%%%%%%
\subsection{LU Factorization With Pivoting}
\label{subsec:lu-with-pivoting}
%%%%%%%%%%%%%%%%%%%%%%%%%%%%%%%%%%%%%%%%%%%%%%%%%%%%%%%%%%%%

With row pivoting, the factorization becomes

\[
\boxed{
PA
=
LU,
}
\]

where \(P\) is a permutation matrix.

To solve

\[
A\mathbf{x}
=
\mathbf{b},
\]

rewrite as

\[
LU\mathbf{x}
=
P\mathbf{b}.
\]

Then perform triangular solves.

%%%%%%%%%%%%%%%%%%%%%%%%%%%%%%%%%%%%%%%%%%%%%%%%%%%%%%%%%%%%
\subsection{Complete Pivoting}
\label{subsec:complete-pivoting}
%%%%%%%%%%%%%%%%%%%%%%%%%%%%%%%%%%%%%%%%%%%%%%%%%%%%%%%%%%%%

Complete pivoting searches the remaining submatrix for the largest
magnitude entry and may interchange both rows and columns.

It can provide stronger pivot control but requires more search and data
movement.

Partial pivoting is much more common in general-purpose dense linear
solvers.

%%%%%%%%%%%%%%%%%%%%%%%%%%%%%%%%%%%%%%%%%%%%%%%%%%%%%%%%%%%%
\subsection{Growth Factor}
\label{subsec:growth-factor}
%%%%%%%%%%%%%%%%%%%%%%%%%%%%%%%%%%%%%%%%%%%%%%%%%%%%%%%%%%%%

During Gaussian elimination, matrix entries may become larger than those in
the original matrix.

A growth factor measures this amplification.

Large growth can amplify rounding effects.

Partial pivoting generally controls growth well in practice, although
specially constructed matrices can produce large worst-case growth.

%%%%%%%%%%%%%%%%%%%%%%%%%%%%%%%%%%%%%%%%%%%%%%%%%%%%%%%%%%%%
\subsection{Computational Cost of Gaussian Elimination}
\label{subsec:gaussian-elimination-cost}
%%%%%%%%%%%%%%%%%%%%%%%%%%%%%%%%%%%%%%%%%%%%%%%%%%%%%%%%%%%%

For a dense \(n\times n\) matrix, LU factorization requires approximately

\[
\boxed{
\frac{2}{3}n^3
}
\]

floating-point arithmetic operations to leading order.

A triangular solve requires only

\[
\boxed{
O(n^2)
}
\]

operations.

Thus factorization dominates the cost.

%%%%%%%%%%%%%%%%%%%%%%%%%%%%%%%%%%%%%%%%%%%%%%%%%%%%%%%%%%%%
\subsection{Multiple Right-Hand Sides}
\label{subsec:multiple-right-hand-sides}
%%%%%%%%%%%%%%%%%%%%%%%%%%%%%%%%%%%%%%%%%%%%%%%%%%%%%%%%%%%%

Suppose

\[
A\mathbf{x}^{(k)}
=
\mathbf{b}^{(k)}
\]

must be solved for many different right-hand sides.

Compute

\[
PA=LU
\]

once.

Then reuse the same factors for every

\[
\mathbf{b}^{(k)}.
\]

This makes factorization especially valuable when the coefficient matrix
does not change.

%%%%%%%%%%%%%%%%%%%%%%%%%%%%%%%%%%%%%%%%%%%%%%%%%%%%%%%%%%%%
\subsection{Symmetric Positive-Definite Matrices}
\label{subsec:symmetric-positive-definite}
%%%%%%%%%%%%%%%%%%%%%%%%%%%%%%%%%%%%%%%%%%%%%%%%%%%%%%%%%%%%

A real matrix \(A\) is symmetric positive definite if

\[
\boxed{
A=A^T
}
\]

and

\[
\boxed{
\mathbf{x}^TA\mathbf{x}>0
}
\]

for every nonzero \(\mathbf{x}\).

Such matrices arise frequently in:

\[
\boxed{
\text{least squares},
}
\]

\[
\boxed{
\text{optimization},
}
\]

and

\[
\boxed{
\text{elliptic PDE discretization}.
}
\]

%%%%%%%%%%%%%%%%%%%%%%%%%%%%%%%%%%%%%%%%%%%%%%%%%%%%%%%%%%%%
\subsection{Cholesky Factorization}
\label{subsec:cholesky-factorization}
%%%%%%%%%%%%%%%%%%%%%%%%%%%%%%%%%%%%%%%%%%%%%%%%%%%%%%%%%%%%

If \(A\) is symmetric positive definite, then

\[
\boxed{
A
=
LL^T,
}
\]

where \(L\) is lower triangular with positive diagonal entries.

This is the Cholesky factorization.

It is more efficient than general LU factorization and preserves symmetry.

%%%%%%%%%%%%%%%%%%%%%%%%%%%%%%%%%%%%%%%%%%%%%%%%%%%%%%%%%%%%
\subsection{Worked Example: Cholesky Factorization}
\label{subsec:worked-cholesky}
%%%%%%%%%%%%%%%%%%%%%%%%%%%%%%%%%%%%%%%%%%%%%%%%%%%%%%%%%%%%

\begin{workedexamplebox}{Factoring a Positive-Definite Matrix}

Let

\[
A
=
\begin{pmatrix}
4&2\\
2&3
\end{pmatrix}.
\]

Assume

\[
L
=
\begin{pmatrix}
\ell_{11}&0\\
\ell_{21}&\ell_{22}
\end{pmatrix}.
\]

From

\[
LL^T=A,
\]

we obtain

\[
\ell_{11}^2=4,
\]

so

\[
\ell_{11}=2.
\]

Next,

\[
\ell_{11}\ell_{21}=2,
\]

so

\[
\ell_{21}=1.
\]

Finally,

\[
\ell_{21}^2+\ell_{22}^2=3,
\]

so

\[
\ell_{22}=\sqrt{2}.
\]

Therefore,

\begin{answerbox}
\[
\boxed{
L
=
\begin{pmatrix}
2&0\\
1&\sqrt{2}
\end{pmatrix}.
}
\]
\end{answerbox}

\end{workedexamplebox}

%%%%%%%%%%%%%%%%%%%%%%%%%%%%%%%%%%%%%%%%%%%%%%%%%%%%%%%%%%%%
\subsection{Cholesky Cost}
\label{subsec:cholesky-cost}
%%%%%%%%%%%%%%%%%%%%%%%%%%%%%%%%%%%%%%%%%%%%%%%%%%%%%%%%%%%%

Dense Cholesky factorization requires approximately

\[
\boxed{
\frac{1}{3}n^3
}
\]

floating-point operations to leading order.

This is roughly half the leading arithmetic cost of general LU
factorization.

%%%%%%%%%%%%%%%%%%%%%%%%%%%%%%%%%%%%%%%%%%%%%%%%%%%%%%%%%%%%
\subsection{Orthogonal Matrices}
\label{subsec:numerical-orthogonal-matrices}
%%%%%%%%%%%%%%%%%%%%%%%%%%%%%%%%%%%%%%%%%%%%%%%%%%%%%%%%%%%%

A square matrix \(Q\) is orthogonal if

\[
\boxed{
Q^TQ=I.
}
\]

Then

\[
\boxed{
Q^{-1}=Q^T.
}
\]

Orthogonal transformations preserve Euclidean norms:

\[
\boxed{
\|Q\mathbf{x}\|_2
=
\|\mathbf{x}\|_2.
}
\]

This makes them numerically valuable.

%%%%%%%%%%%%%%%%%%%%%%%%%%%%%%%%%%%%%%%%%%%%%%%%%%%%%%%%%%%%
\subsection{QR Factorization}
\label{subsec:qr-factorization}
%%%%%%%%%%%%%%%%%%%%%%%%%%%%%%%%%%%%%%%%%%%%%%%%%%%%%%%%%%%%

For

\[
A\in\mathbb{R}^{m\times n},
\qquad
m\geq n,
\]

with linearly independent columns, QR factorization writes

\[
\boxed{
A
=
QR,
}
\]

where \(Q\) has orthonormal columns and \(R\) is upper triangular.

In a thin QR factorization,

\[
\boxed{
Q\in\mathbb{R}^{m\times n},
\qquad
R\in\mathbb{R}^{n\times n}.
}
\]

%%%%%%%%%%%%%%%%%%%%%%%%%%%%%%%%%%%%%%%%%%%%%%%%%%%%%%%%%%%%
\subsection{Gram--Schmidt Orthogonalization}
\label{subsec:gram-schmidt}
%%%%%%%%%%%%%%%%%%%%%%%%%%%%%%%%%%%%%%%%%%%%%%%%%%%%%%%%%%%%

Suppose the columns of \(A\) are

\[
\mathbf{a}_1,\ldots,\mathbf{a}_n.
\]

Gram--Schmidt constructs orthonormal vectors

\[
\mathbf{q}_1,\ldots,\mathbf{q}_n.
\]

Begin with

\[
\boxed{
\mathbf{q}_1
=
\frac{
\mathbf{a}_1
}{
\|\mathbf{a}_1\|_2
}.
}
\]

Then remove from each subsequent column its projections onto previously
constructed orthonormal directions.

%%%%%%%%%%%%%%%%%%%%%%%%%%%%%%%%%%%%%%%%%%%%%%%%%%%%%%%%%%%%
\subsection{Classical Gram--Schmidt}
\label{subsec:classical-gram-schmidt}
%%%%%%%%%%%%%%%%%%%%%%%%%%%%%%%%%%%%%%%%%%%%%%%%%%%%%%%%%%%%

Classical Gram--Schmidt forms

\[
\boxed{
\mathbf{v}_j
=
\mathbf{a}_j
-
\sum_{i=1}^{j-1}
(
\mathbf{q}_i^T\mathbf{a}_j
)
\mathbf{q}_i.
}
\]

Then

\[
\boxed{
\mathbf{q}_j
=
\frac{
\mathbf{v}_j
}{
\|\mathbf{v}_j\|_2
}.
}
\]

In floating-point arithmetic, classical Gram--Schmidt can lose
orthogonality for poorly conditioned column sets.

%%%%%%%%%%%%%%%%%%%%%%%%%%%%%%%%%%%%%%%%%%%%%%%%%%%%%%%%%%%%
\subsection{Modified Gram--Schmidt}
\label{subsec:modified-gram-schmidt}
%%%%%%%%%%%%%%%%%%%%%%%%%%%%%%%%%%%%%%%%%%%%%%%%%%%%%%%%%%%%

Modified Gram--Schmidt removes one projection at a time:

\[
\mathbf{v}
\leftarrow
\mathbf{a}_j.
\]

Then for \(i=1,\ldots,j-1\),

\[
r_{ij}
=
\mathbf{q}_i^T\mathbf{v},
\]

\[
\boxed{
\mathbf{v}
\leftarrow
\mathbf{v}
-
r_{ij}\mathbf{q}_i.
}
\]

This organization generally has better numerical orthogonality properties
than classical Gram--Schmidt.

%%%%%%%%%%%%%%%%%%%%%%%%%%%%%%%%%%%%%%%%%%%%%%%%%%%%%%%%%%%%
\subsection{Householder Reflections}
\label{subsec:householder-reflections}
%%%%%%%%%%%%%%%%%%%%%%%%%%%%%%%%%%%%%%%%%%%%%%%%%%%%%%%%%%%%

A Householder reflector has form

\[
\boxed{
H
=
I
-
2
\frac{
\mathbf{v}\mathbf{v}^T
}{
\mathbf{v}^T\mathbf{v}
}.
}
\]

It is orthogonal and symmetric:

\[
\boxed{
H^T=H,
}
\]

and

\[
\boxed{
H^TH=I.
}
\]

Householder transformations can zero several matrix entries at once.

%%%%%%%%%%%%%%%%%%%%%%%%%%%%%%%%%%%%%%%%%%%%%%%%%%%%%%%%%%%%
\subsection{QR by Householder Transformations}
\label{subsec:qr-householder}
%%%%%%%%%%%%%%%%%%%%%%%%%%%%%%%%%%%%%%%%%%%%%%%%%%%%%%%%%%%%

Householder QR repeatedly constructs reflectors

\[
H_1,H_2,\ldots
\]

that eliminate subdiagonal matrix entries.

The product

\[
\boxed{
Q^T
=
H_k\cdots H_2H_1
}
\]

transforms \(A\) into upper-triangular form:

\[
\boxed{
Q^TA
=
R.
}
\]

Therefore,

\[
\boxed{
A=QR.
}
\]

Householder QR is a standard stable choice for dense least-squares
problems.

%%%%%%%%%%%%%%%%%%%%%%%%%%%%%%%%%%%%%%%%%%%%%%%%%%%%%%%%%%%%
\subsection{Givens Rotations}
\label{subsec:givens-rotations}
%%%%%%%%%%%%%%%%%%%%%%%%%%%%%%%%%%%%%%%%%%%%%%%%%%%%%%%%%%%%

A Givens rotation acts on two coordinates:

\[
\boxed{
G
=
\begin{pmatrix}
c&s\\
-s&c
\end{pmatrix},
}
\]

where

\[
\boxed{
c^2+s^2=1.
}
\]

Parameters \(c\) and \(s\) are chosen to eliminate one selected component
of a vector.

%%%%%%%%%%%%%%%%%%%%%%%%%%%%%%%%%%%%%%%%%%%%%%%%%%%%%%%%%%%%
\subsection{When Givens Rotations Are Useful}
\label{subsec:givens-use}
%%%%%%%%%%%%%%%%%%%%%%%%%%%%%%%%%%%%%%%%%%%%%%%%%%%%%%%%%%%%

Givens rotations are useful when:

\[
\boxed{
\text{only a few entries need to be eliminated},
}
\]

or when

\[
\boxed{
\text{matrix structure should be updated incrementally}.
}
\]

They are especially useful in sparse or updating algorithms.

%%%%%%%%%%%%%%%%%%%%%%%%%%%%%%%%%%%%%%%%%%%%%%%%%%%%%%%%%%%%
\subsection{Least Squares}
\label{subsec:nla-least-squares}
%%%%%%%%%%%%%%%%%%%%%%%%%%%%%%%%%%%%%%%%%%%%%%%%%%%%%%%%%%%%

For an overdetermined system

\[
A\mathbf{x}
\approx
\mathbf{b},
\]

the least-squares problem is

\[
\boxed{
\min_{\mathbf{x}}
\|A\mathbf{x}-\mathbf{b}\|_2.
}
\]

If \(A\) has full column rank, the solution is unique.

%%%%%%%%%%%%%%%%%%%%%%%%%%%%%%%%%%%%%%%%%%%%%%%%%%%%%%%%%%%%
\subsection{Normal Equations}
\label{subsec:nla-normal-equations}
%%%%%%%%%%%%%%%%%%%%%%%%%%%%%%%%%%%%%%%%%%%%%%%%%%%%%%%%%%%%

The least-squares first-order condition is

\[
\boxed{
A^T
(
A\mathbf{x}-\mathbf{b}
)
=
0.
}
\]

Therefore,

\[
\boxed{
A^TA\mathbf{x}
=
A^T\mathbf{b}.
}
\]

These are the normal equations.

%%%%%%%%%%%%%%%%%%%%%%%%%%%%%%%%%%%%%%%%%%%%%%%%%%%%%%%%%%%%
\subsection{Conditioning of Normal Equations}
\label{subsec:normal-equations-conditioning}
%%%%%%%%%%%%%%%%%%%%%%%%%%%%%%%%%%%%%%%%%%%%%%%%%%%%%%%%%%%%

For full-column-rank \(A\),

\[
\boxed{
\kappa_2(A^TA)
=
\kappa_2(A)^2.
}
\]

Thus forming \(A^TA\) can square the condition number.

A moderately ill-conditioned matrix can become much more difficult
numerically.

%%%%%%%%%%%%%%%%%%%%%%%%%%%%%%%%%%%%%%%%%%%%%%%%%%%%%%%%%%%%
\subsection{Least Squares Using QR}
\label{subsec:least-squares-qr}
%%%%%%%%%%%%%%%%%%%%%%%%%%%%%%%%%%%%%%%%%%%%%%%%%%%%%%%%%%%%

Let

\[
A=QR,
\]

with \(Q\) having orthonormal columns.

Then

\[
\|A\mathbf{x}-\mathbf{b}\|_2
=
\|QR\mathbf{x}-\mathbf{b}\|_2.
\]

Using a full orthogonal extension of \(Q\), the least-squares problem
reduces to solving

\[
\boxed{
R\mathbf{x}
=
Q^T\mathbf{b}
}
\]

for the relevant leading components.

This avoids explicitly forming \(A^TA\).

%%%%%%%%%%%%%%%%%%%%%%%%%%%%%%%%%%%%%%%%%%%%%%%%%%%%%%%%%%%%
\subsection{Worked Example: Least Squares by QR}
\label{subsec:worked-least-squares-qr}
%%%%%%%%%%%%%%%%%%%%%%%%%%%%%%%%%%%%%%%%%%%%%%%%%%%%%%%%%%%%

\begin{workedexamplebox}{Using an Existing QR Factorization}

Suppose

\[
A
=
QR,
\]

with

\[
R
=
\begin{pmatrix}
2&1\\
0&3
\end{pmatrix},
\]

and suppose

\[
Q^T\mathbf{b}
=
\begin{pmatrix}
5\\
6
\end{pmatrix}.
\]

Then solve

\[
R\mathbf{x}
=
Q^T\mathbf{b}.
\]

The second equation gives

\[
3x_2=6,
\]

so

\[
x_2=2.
\]

The first equation gives

\[
2x_1+x_2=5.
\]

Thus

\[
2x_1=3,
\]

so

\[
x_1=\frac32.
\]

\begin{answerbox}
\[
\boxed{
\mathbf{x}
=
\begin{pmatrix}
3/2\\
2
\end{pmatrix}.
}
\]
\end{answerbox}

\end{workedexamplebox}

%%%%%%%%%%%%%%%%%%%%%%%%%%%%%%%%%%%%%%%%%%%%%%%%%%%%%%%%%%%%
\subsection{Singular Value Decomposition}
\label{subsec:singular-value-decomposition}
%%%%%%%%%%%%%%%%%%%%%%%%%%%%%%%%%%%%%%%%%%%%%%%%%%%%%%%%%%%%

Every matrix

\[
A\in\mathbb{R}^{m\times n}
\]

has a singular value decomposition

\[
\boxed{
A
=
U\Sigma V^T,
}
\]

where \(U\) and \(V\) are orthogonal matrices and \(\Sigma\) is rectangular
diagonal.

The diagonal entries are the singular values:

\[
\boxed{
\sigma_1
\geq
\sigma_2
\geq
\cdots
\geq
0.
}
\]

%%%%%%%%%%%%%%%%%%%%%%%%%%%%%%%%%%%%%%%%%%%%%%%%%%%%%%%%%%%%
\subsection{Singular Values}
\label{subsec:singular-values}
%%%%%%%%%%%%%%%%%%%%%%%%%%%%%%%%%%%%%%%%%%%%%%%%%%%%%%%%%%%%

The nonzero singular values satisfy

\[
\boxed{
\sigma_i
=
\sqrt{
\lambda_i(A^TA)
}.
}
\]

Thus singular values measure how strongly the matrix stretches vectors in
special orthogonal directions.

%%%%%%%%%%%%%%%%%%%%%%%%%%%%%%%%%%%%%%%%%%%%%%%%%%%%%%%%%%%%
\subsection{Geometric Interpretation of the SVD}
\label{subsec:svd-geometric}
%%%%%%%%%%%%%%%%%%%%%%%%%%%%%%%%%%%%%%%%%%%%%%%%%%%%%%%%%%%%

The transformation

\[
\mathbf{x}
\mapsto
A\mathbf{x}
\]

can be interpreted as:

\begin{enumerate}
    \item rotate or reflect using \(V^T\);
    \item scale coordinate directions using \(\Sigma\);
    \item rotate or reflect using \(U\).
\end{enumerate}

Thus the SVD separates direction and scale.

%%%%%%%%%%%%%%%%%%%%%%%%%%%%%%%%%%%%%%%%%%%%%%%%%%%%%%%%%%%%
\subsection{Rank From the SVD}
\label{subsec:svd-rank}
%%%%%%%%%%%%%%%%%%%%%%%%%%%%%%%%%%%%%%%%%%%%%%%%%%%%%%%%%%%%

The exact rank of \(A\) equals the number of nonzero singular values:

\[
\boxed{
\operatorname{rank}(A)
=
\#\{
i:
\sigma_i>0
\}.
}
\]

In floating-point computation, very small singular values may be treated as
numerically zero.

This leads to the concept of numerical rank.

%%%%%%%%%%%%%%%%%%%%%%%%%%%%%%%%%%%%%%%%%%%%%%%%%%%%%%%%%%%%
\subsection{Numerical Rank}
\label{subsec:numerical-rank}
%%%%%%%%%%%%%%%%%%%%%%%%%%%%%%%%%%%%%%%%%%%%%%%%%%%%%%%%%%%%

A matrix may be mathematically full rank but effectively rank deficient at
the available precision.

A practical numerical rank may be defined by thresholding:

\[
\boxed{
\sigma_i
>
\tau.
}
\]

The tolerance \(\tau\) should reflect matrix scale, dimension, precision,
and the application.

%%%%%%%%%%%%%%%%%%%%%%%%%%%%%%%%%%%%%%%%%%%%%%%%%%%%%%%%%%%%
\subsection{Pseudoinverse}
\label{subsec:moore-penrose-pseudoinverse}
%%%%%%%%%%%%%%%%%%%%%%%%%%%%%%%%%%%%%%%%%%%%%%%%%%%%%%%%%%%%

If

\[
A
=
U\Sigma V^T,
\]

the Moore--Penrose pseudoinverse is

\[
\boxed{
A^+
=
V\Sigma^+U^T,
}
\]

where each nonzero singular value is replaced by its reciprocal.

For suitable least-squares problems,

\[
\boxed{
\mathbf{x}^*
=
A^+\mathbf{b}.
}
\]

%%%%%%%%%%%%%%%%%%%%%%%%%%%%%%%%%%%%%%%%%%%%%%%%%%%%%%%%%%%%
\subsection{Minimum-Norm Solutions}
\label{subsec:minimum-norm-solutions}
%%%%%%%%%%%%%%%%%%%%%%%%%%%%%%%%%%%%%%%%%%%%%%%%%%%%%%%%%%%%

If a linear system has infinitely many solutions, the pseudoinverse selects
the solution with minimum Euclidean norm.

That is,

\[
\boxed{
\mathbf{x}^*
=
A^+\mathbf{b}
}
\]

minimizes

\[
\|\mathbf{x}\|_2
\]

among all exact solutions when the system is consistent.

%%%%%%%%%%%%%%%%%%%%%%%%%%%%%%%%%%%%%%%%%%%%%%%%%%%%%%%%%%%%
\subsection{Low-Rank Approximation}
\label{subsec:low-rank-approximation}
%%%%%%%%%%%%%%%%%%%%%%%%%%%%%%%%%%%%%%%%%%%%%%%%%%%%%%%%%%%%

Suppose

\[
A
=
\sum_{i=1}^{r}
\sigma_i
\mathbf{u}_i
\mathbf{v}_i^T.
\]

A rank-\(k\) approximation is

\[
\boxed{
A_k
=
\sum_{i=1}^{k}
\sigma_i
\mathbf{u}_i
\mathbf{v}_i^T.
}
\]

Only the largest \(k\) singular components are retained.

%%%%%%%%%%%%%%%%%%%%%%%%%%%%%%%%%%%%%%%%%%%%%%%%%%%%%%%%%%%%
\subsection{Eckart--Young Theorem}
\label{subsec:eckart-young}
%%%%%%%%%%%%%%%%%%%%%%%%%%%%%%%%%%%%%%%%%%%%%%%%%%%%%%%%%%%%

\begin{theorem}[Eckart--Young]
Among matrices \(B\) of rank at most \(k\), the truncated SVD \(A_k\)
provides a best approximation to \(A\) in both the spectral and Frobenius
norms.
\end{theorem}

In particular,

\[
\boxed{
\|A-A_k\|_2
=
\sigma_{k+1}.
}
\]

Also,

\[
\boxed{
\|A-A_k\|_F^2
=
\sum_{i=k+1}^{r}
\sigma_i^2.
}
\]

%%%%%%%%%%%%%%%%%%%%%%%%%%%%%%%%%%%%%%%%%%%%%%%%%%%%%%%%%%%%
\subsection{Applications of Low-Rank Approximation}
\label{subsec:low-rank-applications}
%%%%%%%%%%%%%%%%%%%%%%%%%%%%%%%%%%%%%%%%%%%%%%%%%%%%%%%%%%%%

Low-rank approximation is used in:

\[
\boxed{
\text{data compression},
}
\]

\[
\boxed{
\text{noise reduction},
}
\]

\[
\boxed{
\text{latent-factor models},
}
\]

and

\[
\boxed{
\text{dimensionality reduction}.
}
\]

It is a fundamental bridge between numerical linear algebra and data
analysis.

%%%%%%%%%%%%%%%%%%%%%%%%%%%%%%%%%%%%%%%%%%%%%%%%%%%%%%%%%%%%
\subsection{Eigenvalue Problems}
\label{subsec:numerical-eigenvalue-problems}
%%%%%%%%%%%%%%%%%%%%%%%%%%%%%%%%%%%%%%%%%%%%%%%%%%%%%%%%%%%%

An eigenvalue problem seeks

\[
\boxed{
A\mathbf{v}
=
\lambda\mathbf{v},
}
\]

with

\[
\mathbf{v}\neq0.
\]

Numerically, directly computing roots of the characteristic polynomial

\[
\det(A-\lambda I)=0
\]

is generally not the preferred method.

Stable matrix algorithms work directly with \(A\).

%%%%%%%%%%%%%%%%%%%%%%%%%%%%%%%%%%%%%%%%%%%%%%%%%%%%%%%%%%%%
\subsection{Rayleigh Quotient}
\label{subsec:rayleigh-quotient}
%%%%%%%%%%%%%%%%%%%%%%%%%%%%%%%%%%%%%%%%%%%%%%%%%%%%%%%%%%%%

For nonzero \(\mathbf{x}\),

\[
\boxed{
\rho_A(\mathbf{x})
=
\frac{
\mathbf{x}^TA\mathbf{x}
}{
\mathbf{x}^T\mathbf{x}
}.
}
\]

If \(\mathbf{x}\) is an eigenvector of symmetric \(A\), then

\[
\boxed{
\rho_A(\mathbf{x})
=
\lambda.
}
\]

The Rayleigh quotient provides a natural eigenvalue estimate.

%%%%%%%%%%%%%%%%%%%%%%%%%%%%%%%%%%%%%%%%%%%%%%%%%%%%%%%%%%%%
\subsection{Power Method}
\label{subsec:power-method}
%%%%%%%%%%%%%%%%%%%%%%%%%%%%%%%%%%%%%%%%%%%%%%%%%%%%%%%%%%%%

Starting with nonzero vector \(\mathbf{x}_0\), repeatedly compute

\[
\boxed{
\mathbf{y}_{k+1}
=
A\mathbf{x}_k,
}
\]

then normalize:

\[
\boxed{
\mathbf{x}_{k+1}
=
\frac{
\mathbf{y}_{k+1}
}{
\|\mathbf{y}_{k+1}\|
}.
}
\]

Under suitable conditions, \(\mathbf{x}_k\) converges toward an eigenvector
associated with the dominant eigenvalue in magnitude.

%%%%%%%%%%%%%%%%%%%%%%%%%%%%%%%%%%%%%%%%%%%%%%%%%%%%%%%%%%%%
\subsection{Why the Power Method Works}
\label{subsec:why-power-method-works}
%%%%%%%%%%%%%%%%%%%%%%%%%%%%%%%%%%%%%%%%%%%%%%%%%%%%%%%%%%%%

Suppose \(A\) has eigenvectors \(\mathbf{v}_i\) and eigenvalues satisfying

\[
|\lambda_1|
>
|\lambda_2|
\geq
\cdots.
\]

Write

\[
\mathbf{x}_0
=
c_1\mathbf{v}_1
+
c_2\mathbf{v}_2
+
\cdots.
\]

Then

\[
A^k\mathbf{x}_0
=
c_1\lambda_1^k\mathbf{v}_1
+
c_2\lambda_2^k\mathbf{v}_2
+
\cdots.
\]

Factoring \(\lambda_1^k\),

\[
A^k\mathbf{x}_0
=
\lambda_1^k
\left[
c_1\mathbf{v}_1
+
c_2
\left(
\frac{\lambda_2}{\lambda_1}
\right)^k
\mathbf{v}_2
+
\cdots
\right].
\]

The smaller ratios decay when their magnitudes are below \(1\).

%%%%%%%%%%%%%%%%%%%%%%%%%%%%%%%%%%%%%%%%%%%%%%%%%%%%%%%%%%%%
\subsection{Power-Method Convergence Rate}
\label{subsec:power-method-convergence}
%%%%%%%%%%%%%%%%%%%%%%%%%%%%%%%%%%%%%%%%%%%%%%%%%%%%%%%%%%%%

The asymptotic convergence rate is governed by approximately

\[
\boxed{
\left|
\frac{
\lambda_2
}{
\lambda_1
}
\right|.
}
\]

If

\[
|\lambda_2|
\approx
|\lambda_1|,
\]

convergence can be slow.

A large spectral gap improves convergence.

%%%%%%%%%%%%%%%%%%%%%%%%%%%%%%%%%%%%%%%%%%%%%%%%%%%%%%%%%%%%
\subsection{Worked Example: Power Method}
\label{subsec:worked-power-method}
%%%%%%%%%%%%%%%%%%%%%%%%%%%%%%%%%%%%%%%%%%%%%%%%%%%%%%%%%%%%

\begin{workedexamplebox}{One Power Iteration}

Let

\[
A
=
\begin{pmatrix}
3&0\\
0&1
\end{pmatrix},
\]

and

\[
\mathbf{x}_0
=
\begin{pmatrix}
1\\
1
\end{pmatrix}.
\]

Then

\[
A\mathbf{x}_0
=
\begin{pmatrix}
3\\
1
\end{pmatrix}.
\]

Normalize using the Euclidean norm:

\[
\left\|
\begin{pmatrix}
3\\
1
\end{pmatrix}
\right\|_2
=
\sqrt{10}.
\]

Thus

\[
\mathbf{x}_1
=
\frac1{\sqrt{10}}
\begin{pmatrix}
3\\
1
\end{pmatrix}.
\]

Repeated multiplication increasingly aligns the vector with

\[
\begin{pmatrix}
1\\
0
\end{pmatrix},
\]

the eigenvector corresponding to eigenvalue \(3\).

\begin{answerbox}
\[
\boxed{
\lambda_{\max}=3.
}
\]
\end{answerbox}

\end{workedexamplebox}

%%%%%%%%%%%%%%%%%%%%%%%%%%%%%%%%%%%%%%%%%%%%%%%%%%%%%%%%%%%%
\subsection{Inverse Iteration}
\label{subsec:inverse-iteration}
%%%%%%%%%%%%%%%%%%%%%%%%%%%%%%%%%%%%%%%%%%%%%%%%%%%%%%%%%%%%

Inverse iteration repeatedly solves

\[
\boxed{
A\mathbf{y}_{k+1}
=
\mathbf{x}_k
}
\]

and normalizes \(\mathbf{y}_{k+1}\).

Because eigenvalues of \(A^{-1}\) are

\[
\frac1{\lambda_i},
\]

inverse iteration tends toward an eigenvector associated with the
smallest-magnitude eigenvalue of \(A\), under suitable conditions.

%%%%%%%%%%%%%%%%%%%%%%%%%%%%%%%%%%%%%%%%%%%%%%%%%%%%%%%%%%%%
\subsection{Shifted Inverse Iteration}
\label{subsec:shifted-inverse-iteration}
%%%%%%%%%%%%%%%%%%%%%%%%%%%%%%%%%%%%%%%%%%%%%%%%%%%%%%%%%%%%

To target an eigenvalue near shift \(\mu\), solve repeatedly

\[
\boxed{
(A-\mu I)
\mathbf{y}_{k+1}
=
\mathbf{x}_k.
}
\]

Eigenvalues of

\[
(A-\mu I)^{-1}
\]

are

\[
\frac1{\lambda_i-\mu}.
\]

Thus the eigenvalue nearest \(\mu\) becomes dominant in magnitude.

%%%%%%%%%%%%%%%%%%%%%%%%%%%%%%%%%%%%%%%%%%%%%%%%%%%%%%%%%%%%
\subsection{Rayleigh Quotient Iteration}
\label{subsec:rayleigh-quotient-iteration}
%%%%%%%%%%%%%%%%%%%%%%%%%%%%%%%%%%%%%%%%%%%%%%%%%%%%%%%%%%%%

Rayleigh quotient iteration updates the shift dynamically:

\[
\boxed{
\mu_k
=
\frac{
\mathbf{x}_k^TA\mathbf{x}_k
}{
\mathbf{x}_k^T\mathbf{x}_k
}.
}
\]

Then solve

\[
\boxed{
(A-\mu_kI)
\mathbf{y}_{k+1}
=
\mathbf{x}_k,
}
\]

and normalize.

For symmetric matrices near an eigenvector, convergence can be extremely
rapid.

%%%%%%%%%%%%%%%%%%%%%%%%%%%%%%%%%%%%%%%%%%%%%%%%%%%%%%%%%%%%
\subsection{QR Iteration}
\label{subsec:qr-iteration}
%%%%%%%%%%%%%%%%%%%%%%%%%%%%%%%%%%%%%%%%%%%%%%%%%%%%%%%%%%%%

A basic QR eigenvalue iteration performs

\[
A_k
=
Q_kR_k,
\]

then forms

\[
\boxed{
A_{k+1}
=
R_kQ_k.
}
\]

Because

\[
A_{k+1}
=
Q_k^TA_kQ_k,
\]

the matrices are orthogonally similar and therefore share eigenvalues.

Repeated iteration can drive \(A_k\) toward triangular or nearly diagonal
form.

%%%%%%%%%%%%%%%%%%%%%%%%%%%%%%%%%%%%%%%%%%%%%%%%%%%%%%%%%%%%
\subsection{Shifted QR Iteration}
\label{subsec:shifted-qr-iteration}
%%%%%%%%%%%%%%%%%%%%%%%%%%%%%%%%%%%%%%%%%%%%%%%%%%%%%%%%%%%%

Practical QR algorithms use shifts:

\[
A_k-\mu_kI
=
Q_kR_k,
\]

followed by

\[
\boxed{
A_{k+1}
=
R_kQ_k+\mu_kI.
}
\]

Carefully selected shifts can greatly accelerate convergence.

%%%%%%%%%%%%%%%%%%%%%%%%%%%%%%%%%%%%%%%%%%%%%%%%%%%%%%%%%%%%
\subsection{Hessenberg Reduction}
\label{subsec:hessenberg-reduction}
%%%%%%%%%%%%%%%%%%%%%%%%%%%%%%%%%%%%%%%%%%%%%%%%%%%%%%%%%%%%

Before applying QR iteration to a general dense matrix, it is efficient to
reduce \(A\) to upper Hessenberg form:

\[
\boxed{
h_{ij}=0
\qquad
\text{for}
\qquad
i>j+1.
}
\]

Householder transformations can perform this reduction while preserving
eigenvalues.

%%%%%%%%%%%%%%%%%%%%%%%%%%%%%%%%%%%%%%%%%%%%%%%%%%%%%%%%%%%%
\subsection{Symmetric Tridiagonal Reduction}
\label{subsec:symmetric-tridiagonal-reduction}
%%%%%%%%%%%%%%%%%%%%%%%%%%%%%%%%%%%%%%%%%%%%%%%%%%%%%%%%%%%%

If \(A\) is symmetric, orthogonal similarity transformations can reduce it
to a symmetric tridiagonal matrix.

This preserves symmetry and significantly reduces the cost of later
eigenvalue iterations.

%%%%%%%%%%%%%%%%%%%%%%%%%%%%%%%%%%%%%%%%%%%%%%%%%%%%%%%%%%%%
\subsection{Direct and Iterative Linear Solvers}
\label{subsec:direct-versus-iterative-linear-solvers}
%%%%%%%%%%%%%%%%%%%%%%%%%%%%%%%%%%%%%%%%%%%%%%%%%%%%%%%%%%%%

Direct methods such as LU attempt to obtain the solution after a finite
sequence of arithmetic operations, apart from floating-point effects.

Iterative methods generate approximations:

\[
\boxed{
\mathbf{x}^{(0)},
\mathbf{x}^{(1)},
\mathbf{x}^{(2)},
\ldots
}
\]

and stop when the approximation is sufficiently accurate.

%%%%%%%%%%%%%%%%%%%%%%%%%%%%%%%%%%%%%%%%%%%%%%%%%%%%%%%%%%%%
\subsection{Why Use Iterative Methods?}
\label{subsec:why-iterative-linear-solvers}
%%%%%%%%%%%%%%%%%%%%%%%%%%%%%%%%%%%%%%%%%%%%%%%%%%%%%%%%%%%%

Iterative methods can be preferable when:

\[
\boxed{
A
\text{ is very large},
}
\]

\[
\boxed{
A
\text{ is sparse},
}
\]

or

\[
\boxed{
\text{only moderate solution accuracy is needed}.
}
\]

They may avoid the memory and fill-in costs of direct factorization.

%%%%%%%%%%%%%%%%%%%%%%%%%%%%%%%%%%%%%%%%%%%%%%%%%%%%%%%%%%%%
\subsection{Matrix Splitting}
\label{subsec:matrix-splitting}
%%%%%%%%%%%%%%%%%%%%%%%%%%%%%%%%%%%%%%%%%%%%%%%%%%%%%%%%%%%%

Write

\[
\boxed{
A
=
M-N,
}
\]

where \(M\) is easy to invert.

Then

\[
A\mathbf{x}
=
\mathbf{b}
\]

becomes

\[
M\mathbf{x}
=
N\mathbf{x}
+
\mathbf{b}.
\]

This motivates iteration

\[
\boxed{
M\mathbf{x}^{(k+1)}
=
N\mathbf{x}^{(k)}
+
\mathbf{b}.
}
\]

Thus,

\[
\boxed{
\mathbf{x}^{(k+1)}
=
M^{-1}N\mathbf{x}^{(k)}
+
M^{-1}\mathbf{b}.
}
\]

%%%%%%%%%%%%%%%%%%%%%%%%%%%%%%%%%%%%%%%%%%%%%%%%%%%%%%%%%%%%
\subsection{Iteration Matrix}
\label{subsec:iteration-matrix}
%%%%%%%%%%%%%%%%%%%%%%%%%%%%%%%%%%%%%%%%%%%%%%%%%%%%%%%%%%%%

Define

\[
\boxed{
G
=
M^{-1}N.
}
\]

Then

\[
\boxed{
\mathbf{x}^{(k+1)}
=
G\mathbf{x}^{(k)}
+
\mathbf{c}.
}
\]

The error satisfies

\[
\boxed{
\mathbf{e}^{(k+1)}
=
G\mathbf{e}^{(k)}.
}
\]

Therefore,

\[
\boxed{
\mathbf{e}^{(k)}
=
G^k
\mathbf{e}^{(0)}.
}
\]

%%%%%%%%%%%%%%%%%%%%%%%%%%%%%%%%%%%%%%%%%%%%%%%%%%%%%%%%%%%%
\subsection{Spectral Radius}
\label{subsec:spectral-radius-iteration}
%%%%%%%%%%%%%%%%%%%%%%%%%%%%%%%%%%%%%%%%%%%%%%%%%%%%%%%%%%%%

The spectral radius of \(G\) is

\[
\boxed{
\rho(G)
=
\max_i
|\lambda_i(G)|.
}
\]

For a stationary linear iteration, convergence for every initial guess
occurs precisely when

\[
\boxed{
\rho(G)<1.
}
\]

This provides a central convergence criterion.

%%%%%%%%%%%%%%%%%%%%%%%%%%%%%%%%%%%%%%%%%%%%%%%%%%%%%%%%%%%%
\subsection{Jacobi Iteration}
\label{subsec:jacobi-iteration}
%%%%%%%%%%%%%%%%%%%%%%%%%%%%%%%%%%%%%%%%%%%%%%%%%%%%%%%%%%%%

Write

\[
A
=
D+L+U,
\]

where:

\[
D
=
\text{diagonal part},
\]

\[
L
=
\text{strictly lower triangular part},
\]

and

\[
U
=
\text{strictly upper triangular part}.
\]

Jacobi iteration is

\[
\boxed{
D\mathbf{x}^{(k+1)}
=
\mathbf{b}
-
(L+U)
\mathbf{x}^{(k)}.
}
\]

Thus,

\[
\boxed{
\mathbf{x}^{(k+1)}
=
D^{-1}
\left[
\mathbf{b}
-
(L+U)\mathbf{x}^{(k)}
\right].
}
\]

%%%%%%%%%%%%%%%%%%%%%%%%%%%%%%%%%%%%%%%%%%%%%%%%%%%%%%%%%%%%
\subsection{Component Form of Jacobi}
\label{subsec:jacobi-component-form}
%%%%%%%%%%%%%%%%%%%%%%%%%%%%%%%%%%%%%%%%%%%%%%%%%%%%%%%%%%%%

For each component,

\[
\boxed{
x_i^{(k+1)}
=
\frac1{a_{ii}}
\left[
b_i
-
\sum_{j\neq i}
a_{ij}
x_j^{(k)}
\right].
}
\]

Every updated component uses only values from the previous iteration.

This makes Jacobi naturally parallelizable.

%%%%%%%%%%%%%%%%%%%%%%%%%%%%%%%%%%%%%%%%%%%%%%%%%%%%%%%%%%%%
\subsection{Gauss--Seidel Iteration}
\label{subsec:gauss-seidel}
%%%%%%%%%%%%%%%%%%%%%%%%%%%%%%%%%%%%%%%%%%%%%%%%%%%%%%%%%%%%

Gauss--Seidel uses newly updated values immediately.

Its matrix form is

\[
\boxed{
(D+L)
\mathbf{x}^{(k+1)}
=
\mathbf{b}
-
U\mathbf{x}^{(k)}.
}
\]

Therefore,

\[
\boxed{
\mathbf{x}^{(k+1)}
=
-(D+L)^{-1}
U\mathbf{x}^{(k)}
+
(D+L)^{-1}\mathbf{b}.
}
\]

%%%%%%%%%%%%%%%%%%%%%%%%%%%%%%%%%%%%%%%%%%%%%%%%%%%%%%%%%%%%
\subsection{Component Form of Gauss--Seidel}
\label{subsec:gauss-seidel-component}
%%%%%%%%%%%%%%%%%%%%%%%%%%%%%%%%%%%%%%%%%%%%%%%%%%%%%%%%%%%%

For each \(i\),

\[
\boxed{
x_i^{(k+1)}
=
\frac1{a_{ii}}
\left[
b_i
-
\sum_{j<i}
a_{ij}
x_j^{(k+1)}
-
\sum_{j>i}
a_{ij}
x_j^{(k)}
\right].
}
\]

Updated values are used as soon as they become available.

%%%%%%%%%%%%%%%%%%%%%%%%%%%%%%%%%%%%%%%%%%%%%%%%%%%%%%%%%%%%
\subsection{Worked Example: Jacobi Iteration}
\label{subsec:worked-jacobi}
%%%%%%%%%%%%%%%%%%%%%%%%%%%%%%%%%%%%%%%%%%%%%%%%%%%%%%%%%%%%

\begin{workedexamplebox}{One Jacobi Step}

Consider

\[
\begin{aligned}
4x+y&=9,\\
x+3y&=7.
\end{aligned}
\]

Solve each equation for its diagonal variable:

\[
x
=
\frac{9-y}{4},
\]

\[
y
=
\frac{7-x}{3}.
\]

Starting with

\[
(x^{(0)},y^{(0)})
=
(0,0),
\]

the first Jacobi iteration gives

\[
x^{(1)}
=
\frac94
=
2.25,
\]

and

\[
y^{(1)}
=
\frac73
\approx2.3333.
\]

\begin{answerbox}
\[
\boxed{
\mathbf{x}^{(1)}
\approx
\begin{pmatrix}
2.25\\
2.3333
\end{pmatrix}.
}
\]
\end{answerbox}

\end{workedexamplebox}

%%%%%%%%%%%%%%%%%%%%%%%%%%%%%%%%%%%%%%%%%%%%%%%%%%%%%%%%%%%%
\subsection{Diagonal Dominance}
\label{subsec:diagonal-dominance}
%%%%%%%%%%%%%%%%%%%%%%%%%%%%%%%%%%%%%%%%%%%%%%%%%%%%%%%%%%%%

A matrix is strictly row diagonally dominant if

\[
\boxed{
|a_{ii}|
>
\sum_{j\neq i}
|a_{ij}|
}
\]

for every row \(i\).

Strict diagonal dominance provides a useful sufficient condition for
convergence of standard Jacobi and Gauss--Seidel iterations.

It is not a necessary condition.

%%%%%%%%%%%%%%%%%%%%%%%%%%%%%%%%%%%%%%%%%%%%%%%%%%%%%%%%%%%%
\subsection{Successive Over-Relaxation}
\label{subsec:successive-over-relaxation}
%%%%%%%%%%%%%%%%%%%%%%%%%%%%%%%%%%%%%%%%%%%%%%%%%%%%%%%%%%%%

Successive Over-Relaxation, abbreviated SOR, modifies the Gauss--Seidel
step using relaxation parameter

\[
\boxed{
\omega.
}
\]

Conceptually,

\[
\boxed{
\mathbf{x}^{(k+1)}
=
\mathbf{x}^{(k)}
+
\omega
\left(
\mathbf{x}_{GS}^{(k+1)}
-
\mathbf{x}^{(k)}
\right).
}
\]

For

\[
\omega=1,
\]

SOR reduces to Gauss--Seidel.

%%%%%%%%%%%%%%%%%%%%%%%%%%%%%%%%%%%%%%%%%%%%%%%%%%%%%%%%%%%%
\subsection{Under- and Over-Relaxation}
\label{subsec:under-over-relaxation}
%%%%%%%%%%%%%%%%%%%%%%%%%%%%%%%%%%%%%%%%%%%%%%%%%%%%%%%%%%%%

If

\[
0<\omega<1,
\]

the method is under-relaxed.

If

\[
1<\omega<2,
\]

it is over-relaxed.

A good \(\omega\) can accelerate convergence for suitable problems.

A poor choice may slow convergence or cause divergence.

%%%%%%%%%%%%%%%%%%%%%%%%%%%%%%%%%%%%%%%%%%%%%%%%%%%%%%%%%%%%
\subsection{Krylov Subspaces}
\label{subsec:krylov-subspaces}
%%%%%%%%%%%%%%%%%%%%%%%%%%%%%%%%%%%%%%%%%%%%%%%%%%%%%%%%%%%%

Given matrix \(A\) and vector \(\mathbf{r}_0\), the \(k\)-th Krylov
subspace is

\[
\boxed{
\mathcal{K}_k
(
A,\mathbf{r}_0
)
=
\operatorname{span}
\{
\mathbf{r}_0,
A\mathbf{r}_0,
A^2\mathbf{r}_0,
\ldots,
A^{k-1}\mathbf{r}_0
\}.
}
\]

Many modern iterative methods search for corrections within Krylov
subspaces.

%%%%%%%%%%%%%%%%%%%%%%%%%%%%%%%%%%%%%%%%%%%%%%%%%%%%%%%%%%%%
\subsection{Conjugate Gradient Method}
\label{subsec:conjugate-gradient}
%%%%%%%%%%%%%%%%%%%%%%%%%%%%%%%%%%%%%%%%%%%%%%%%%%%%%%%%%%%%

The Conjugate Gradient method, abbreviated CG, solves

\[
\boxed{
A\mathbf{x}
=
\mathbf{b}
}
\]

when \(A\) is symmetric positive definite.

It can also be interpreted as minimizing the quadratic function

\[
\boxed{
\phi(\mathbf{x})
=
\frac12
\mathbf{x}^TA\mathbf{x}
-
\mathbf{b}^T\mathbf{x}.
}
\]

%%%%%%%%%%%%%%%%%%%%%%%%%%%%%%%%%%%%%%%%%%%%%%%%%%%%%%%%%%%%
\subsection{Gradient of the CG Quadratic}
\label{subsec:cg-gradient}
%%%%%%%%%%%%%%%%%%%%%%%%%%%%%%%%%%%%%%%%%%%%%%%%%%%%%%%%%%%%

The gradient is

\[
\boxed{
\nabla\phi(\mathbf{x})
=
A\mathbf{x}
-
\mathbf{b}.
}
\]

Thus the negative gradient is the residual:

\[
\boxed{
-\nabla\phi(\mathbf{x})
=
\mathbf{b}
-
A\mathbf{x}
=
\mathbf{r}.
}
\]

CG chooses search directions more intelligently than ordinary steepest
descent.

%%%%%%%%%%%%%%%%%%%%%%%%%%%%%%%%%%%%%%%%%%%%%%%%%%%%%%%%%%%%
\subsection{\(A\)-Conjugate Directions}
\label{subsec:a-conjugate-directions}
%%%%%%%%%%%%%%%%%%%%%%%%%%%%%%%%%%%%%%%%%%%%%%%%%%%%%%%%%%%%

Vectors \(\mathbf{p}_i\) and \(\mathbf{p}_j\) are \(A\)-conjugate if

\[
\boxed{
\mathbf{p}_i^TA\mathbf{p}_j
=
0.
}
\]

For symmetric positive-definite \(A\), CG constructs mutually
\(A\)-conjugate search directions.

In exact arithmetic, CG terminates in at most \(n\) iterations for an
\(n\times n\) system.

%%%%%%%%%%%%%%%%%%%%%%%%%%%%%%%%%%%%%%%%%%%%%%%%%%%%%%%%%%%%
\subsection{Conjugate Gradient Algorithm}
\label{subsec:cg-algorithm}
%%%%%%%%%%%%%%%%%%%%%%%%%%%%%%%%%%%%%%%%%%%%%%%%%%%%%%%%%%%%

Start from \(\mathbf{x}_0\) and compute

\[
\boxed{
\mathbf{r}_0
=
\mathbf{b}
-
A\mathbf{x}_0,
}
\]

\[
\boxed{
\mathbf{p}_0
=
\mathbf{r}_0.
}
\]

Then iterate:

\[
\boxed{
\alpha_k
=
\frac{
\mathbf{r}_k^T\mathbf{r}_k
}{
\mathbf{p}_k^TA\mathbf{p}_k
},
}
\]

\[
\boxed{
\mathbf{x}_{k+1}
=
\mathbf{x}_k
+
\alpha_k
\mathbf{p}_k,
}
\]

\[
\boxed{
\mathbf{r}_{k+1}
=
\mathbf{r}_k
-
\alpha_k
A\mathbf{p}_k,
}
\]

\[
\boxed{
\beta_k
=
\frac{
\mathbf{r}_{k+1}^T\mathbf{r}_{k+1}
}{
\mathbf{r}_k^T\mathbf{r}_k
},
}
\]

and

\[
\boxed{
\mathbf{p}_{k+1}
=
\mathbf{r}_{k+1}
+
\beta_k
\mathbf{p}_k.
}
\]

%%%%%%%%%%%%%%%%%%%%%%%%%%%%%%%%%%%%%%%%%%%%%%%%%%%%%%%%%%%%
\subsection{CG Error Bound}
\label{subsec:cg-error-bound}
%%%%%%%%%%%%%%%%%%%%%%%%%%%%%%%%%%%%%%%%%%%%%%%%%%%%%%%%%%%%

For symmetric positive-definite \(A\), a classical bound is

\[
\boxed{
\|\mathbf{x}_k-\mathbf{x}^*\|_A
\leq
2
\left(
\frac{
\sqrt{\kappa_2(A)}-1
}{
\sqrt{\kappa_2(A)}+1
}
\right)^k
\|\mathbf{x}_0-\mathbf{x}^*\|_A,
}
\]

where

\[
\boxed{
\|\mathbf{v}\|_A
=
\sqrt{
\mathbf{v}^TA\mathbf{v}
}.
}
\]

This shows why conditioning strongly affects CG convergence.

%%%%%%%%%%%%%%%%%%%%%%%%%%%%%%%%%%%%%%%%%%%%%%%%%%%%%%%%%%%%
\subsection{GMRES}
\label{subsec:gmres}
%%%%%%%%%%%%%%%%%%%%%%%%%%%%%%%%%%%%%%%%%%%%%%%%%%%%%%%%%%%%

The Generalized Minimal Residual method, abbreviated GMRES, is a Krylov
method for nonsymmetric linear systems.

At iteration \(k\), it seeks

\[
\boxed{
\mathbf{x}_k
\in
\mathbf{x}_0
+
\mathcal{K}_k
(
A,\mathbf{r}_0
)
}
\]

that minimizes the residual norm

\[
\boxed{
\|\mathbf{b}-A\mathbf{x}_k\|_2.
}
\]

%%%%%%%%%%%%%%%%%%%%%%%%%%%%%%%%%%%%%%%%%%%%%%%%%%%%%%%%%%%%
\subsection{Arnoldi Process}
\label{subsec:arnoldi-process}
%%%%%%%%%%%%%%%%%%%%%%%%%%%%%%%%%%%%%%%%%%%%%%%%%%%%%%%%%%%%

GMRES constructs an orthonormal basis of the Krylov subspace using the
Arnoldi process.

Arnoldi produces relation

\[
\boxed{
AQ_k
=
Q_{k+1}
\overline{H}_k,
}
\]

where \(\overline{H}_k\) is upper Hessenberg.

The large linear problem is reduced to a smaller least-squares problem.

%%%%%%%%%%%%%%%%%%%%%%%%%%%%%%%%%%%%%%%%%%%%%%%%%%%%%%%%%%%%
\subsection{GMRES Restarting}
\label{subsec:gmres-restarting}
%%%%%%%%%%%%%%%%%%%%%%%%%%%%%%%%%%%%%%%%%%%%%%%%%%%%%%%%%%%%

Full GMRES stores an increasing number of Krylov basis vectors.

For large problems, memory and orthogonalization cost increase with each
iteration.

Restarted GMRES limits the subspace dimension and periodically starts a
new Krylov cycle.

Restarting reduces storage but can slow convergence.

%%%%%%%%%%%%%%%%%%%%%%%%%%%%%%%%%%%%%%%%%%%%%%%%%%%%%%%%%%%%
\subsection{Preconditioning}
\label{subsec:preconditioning}
%%%%%%%%%%%%%%%%%%%%%%%%%%%%%%%%%%%%%%%%%%%%%%%%%%%%%%%%%%%%

A preconditioner is a matrix \(M\) that approximates \(A\) but is easier to
solve with.

Instead of solving

\[
A\mathbf{x}
=
\mathbf{b},
\]

one may solve a transformed system such as

\[
\boxed{
M^{-1}A\mathbf{x}
=
M^{-1}\mathbf{b}.
}
\]

The goal is to improve numerical properties and accelerate iterative
convergence.

%%%%%%%%%%%%%%%%%%%%%%%%%%%%%%%%%%%%%%%%%%%%%%%%%%%%%%%%%%%%
\subsection{What Makes a Good Preconditioner?}
\label{subsec:good-preconditioner}
%%%%%%%%%%%%%%%%%%%%%%%%%%%%%%%%%%%%%%%%%%%%%%%%%%%%%%%%%%%%

A good preconditioner should:

\[
\boxed{
\text{approximate }A\text{ sufficiently well},
}
\]

while also being

\[
\boxed{
\text{cheap to construct}
}
\]

and

\[
\boxed{
\text{cheap to apply}.
}
\]

An excellent approximation that is as expensive as solving the original
system may provide little practical benefit.

%%%%%%%%%%%%%%%%%%%%%%%%%%%%%%%%%%%%%%%%%%%%%%%%%%%%%%%%%%%%
\subsection{Jacobi Preconditioning}
\label{subsec:jacobi-preconditioning}
%%%%%%%%%%%%%%%%%%%%%%%%%%%%%%%%%%%%%%%%%%%%%%%%%%%%%%%%%%%%

A simple preconditioner is

\[
\boxed{
M
=
D,
}
\]

where \(D\) is the diagonal of \(A\).

Then applying \(M^{-1}\) is inexpensive.

This can improve scaling, although it may be insufficient for difficult
systems.

%%%%%%%%%%%%%%%%%%%%%%%%%%%%%%%%%%%%%%%%%%%%%%%%%%%%%%%%%%%%
\subsection{Incomplete Factorization}
\label{subsec:incomplete-factorization}
%%%%%%%%%%%%%%%%%%%%%%%%%%%%%%%%%%%%%%%%%%%%%%%%%%%%%%%%%%%%

An incomplete LU factorization constructs approximate factors

\[
\boxed{
A
\approx
LU
}
\]

while discarding selected fill-in entries.

Similarly, incomplete Cholesky can be used for suitable symmetric
positive-definite systems.

These factorizations are widely used as preconditioners.

%%%%%%%%%%%%%%%%%%%%%%%%%%%%%%%%%%%%%%%%%%%%%%%%%%%%%%%%%%%%
\subsection{Sparse Matrix Storage}
\label{subsec:sparse-matrix-storage}
%%%%%%%%%%%%%%%%%%%%%%%%%%%%%%%%%%%%%%%%%%%%%%%%%%%%%%%%%%%%

Sparse matrices should generally store only nonzero values and structural
information.

Common ideas include storing:

\[
\boxed{
\text{nonzero values},
}
\]

\[
\boxed{
\text{column indices},
}
\]

and

\[
\boxed{
\text{row pointers}.
}
\]

Compressed sparse row and compressed sparse column formats are widely
used.

%%%%%%%%%%%%%%%%%%%%%%%%%%%%%%%%%%%%%%%%%%%%%%%%%%%%%%%%%%%%
\subsection{Sparse Matrix-Vector Multiplication}
\label{subsec:sparse-matrix-vector-product}
%%%%%%%%%%%%%%%%%%%%%%%%%%%%%%%%%%%%%%%%%%%%%%%%%%%%%%%%%%%%

For a sparse matrix with

\[
\operatorname{nnz}(A)
\]

nonzero entries, the matrix-vector product

\[
A\mathbf{x}
\]

can often be computed in approximately

\[
\boxed{
O
(
\operatorname{nnz}(A)
)
}
\]

operations.

This makes matrix-vector-product-based Krylov methods attractive.

%%%%%%%%%%%%%%%%%%%%%%%%%%%%%%%%%%%%%%%%%%%%%%%%%%%%%%%%%%%%
\subsection{Fill-In}
\label{subsec:fill-in}
%%%%%%%%%%%%%%%%%%%%%%%%%%%%%%%%%%%%%%%%%%%%%%%%%%%%%%%%%%%%

During sparse LU or Cholesky factorization, entries that were originally
zero may become nonzero.

This is called fill-in.

Excessive fill-in can cause:

\[
\boxed{
\text{large memory use},
}
\]

and

\[
\boxed{
\text{large computational cost}.
}
\]

%%%%%%%%%%%%%%%%%%%%%%%%%%%%%%%%%%%%%%%%%%%%%%%%%%%%%%%%%%%%
\subsection{Reordering Sparse Systems}
\label{subsec:sparse-reordering}
%%%%%%%%%%%%%%%%%%%%%%%%%%%%%%%%%%%%%%%%%%%%%%%%%%%%%%%%%%%%

Changing the order of rows and columns can dramatically affect fill-in.

Sparse direct solvers therefore use graph-based ordering strategies before
factorization.

The algebraically equivalent system may become much cheaper to solve after
an appropriate permutation.

%%%%%%%%%%%%%%%%%%%%%%%%%%%%%%%%%%%%%%%%%%%%%%%%%%%%%%%%%%%%
\subsection{Iterative Refinement}
\label{subsec:iterative-refinement}
%%%%%%%%%%%%%%%%%%%%%%%%%%%%%%%%%%%%%%%%%%%%%%%%%%%%%%%%%%%%

Suppose an approximate solution \(\widehat{\mathbf{x}}\) is available.

Compute residual

\[
\boxed{
\mathbf{r}
=
\mathbf{b}
-
A\widehat{\mathbf{x}}.
}
\]

Solve the correction equation

\[
\boxed{
A\mathbf{d}
=
\mathbf{r}.
}
\]

Then update

\[
\boxed{
\widehat{\mathbf{x}}
\leftarrow
\widehat{\mathbf{x}}
+
\mathbf{d}.
}
\]

Repeated corrections can improve solution accuracy.

%%%%%%%%%%%%%%%%%%%%%%%%%%%%%%%%%%%%%%%%%%%%%%%%%%%%%%%%%%%%
\subsection{Mixed-Precision Refinement Preview}
\label{subsec:mixed-precision-refinement}
%%%%%%%%%%%%%%%%%%%%%%%%%%%%%%%%%%%%%%%%%%%%%%%%%%%%%%%%%%%%

A factorization may sometimes be computed in lower precision for speed,
while residuals and corrections are evaluated in higher precision.

Under suitable conditioning and algorithmic conditions, iterative
refinement can recover high-accuracy solutions.

This is an important modern use of mixed-precision arithmetic.

%%%%%%%%%%%%%%%%%%%%%%%%%%%%%%%%%%%%%%%%%%%%%%%%%%%%%%%%%%%%
\subsection{Backward Stability of Linear Solvers}
\label{subsec:backward-stability-linear-solvers}
%%%%%%%%%%%%%%%%%%%%%%%%%%%%%%%%%%%%%%%%%%%%%%%%%%%%%%%%%%%%

A computed solution \(\widehat{\mathbf{x}}\) is backward stable if it
solves a nearby problem exactly:

\[
\boxed{
(A+\Delta A)
\widehat{\mathbf{x}}
=
\mathbf{b}+\Delta\mathbf{b},
}
\]

where the perturbations are small relative to the original data.

Gaussian elimination with partial pivoting is generally backward stable in
practice, subject to growth-factor considerations.

%%%%%%%%%%%%%%%%%%%%%%%%%%%%%%%%%%%%%%%%%%%%%%%%%%%%%%%%%%%%
\subsection{Conditioning and Backward Stability Together}
\label{subsec:conditioning-backward-stability}
%%%%%%%%%%%%%%%%%%%%%%%%%%%%%%%%%%%%%%%%%%%%%%%%%%%%%%%%%%%%

A backward-stable method solves a nearby problem accurately.

If the original problem is well conditioned, the nearby exact solution is
also near the original exact solution.

Thus:

\[
\boxed{
\text{backward stability}
+
\text{good conditioning}
\Rightarrow
\text{small forward error}.
}
\]

This principle is central to numerical linear algebra.

%%%%%%%%%%%%%%%%%%%%%%%%%%%%%%%%%%%%%%%%%%%%%%%%%%%%%%%%%%%%
\subsection{Scaling}
\label{subsec:matrix-scaling}
%%%%%%%%%%%%%%%%%%%%%%%%%%%%%%%%%%%%%%%%%%%%%%%%%%%%%%%%%%%%

If rows or columns have very different magnitudes, numerical performance
can suffer.

Scaling replaces the system with an equivalent or closely related one in
which magnitudes are better balanced.

For example,

\[
\boxed{
D_rAD_c\mathbf{y}
=
D_r\mathbf{b},
}
\]

with diagonal scaling matrices \(D_r\) and \(D_c\).

%%%%%%%%%%%%%%%%%%%%%%%%%%%%%%%%%%%%%%%%%%%%%%%%%%%%%%%%%%%%
\subsection{Matrix Equilibration}
\label{subsec:matrix-equilibration}
%%%%%%%%%%%%%%%%%%%%%%%%%%%%%%%%%%%%%%%%%%%%%%%%%%%%%%%%%%%%

Equilibration attempts to scale matrix rows or columns so their norms are
similar.

The goal is not necessarily to minimize the mathematical condition number
exactly.

Rather, it can improve numerical behavior and reduce poor scaling effects.

%%%%%%%%%%%%%%%%%%%%%%%%%%%%%%%%%%%%%%%%%%%%%%%%%%%%%%%%%%%%
\subsection{Rank-Revealing Factorizations}
\label{subsec:rank-revealing-factorizations}
%%%%%%%%%%%%%%%%%%%%%%%%%%%%%%%%%%%%%%%%%%%%%%%%%%%%%%%%%%%%

When a matrix may be nearly rank deficient, one wants a factorization that
exposes weak directions.

Useful approaches include:

\[
\boxed{
\text{SVD},
}
\]

and

\[
\boxed{
\text{pivoted QR}.
}
\]

These methods help estimate numerical rank and solve rank-deficient least-
squares problems.

%%%%%%%%%%%%%%%%%%%%%%%%%%%%%%%%%%%%%%%%%%%%%%%%%%%%%%%%%%%%
\subsection{QR With Column Pivoting}
\label{subsec:qr-column-pivoting}
%%%%%%%%%%%%%%%%%%%%%%%%%%%%%%%%%%%%%%%%%%%%%%%%%%%%%%%%%%%%

A pivoted QR factorization has form

\[
\boxed{
AP
=
QR,
}
\]

where \(P\) permutes columns.

The permutation attempts to place more linearly independent or larger
directions earlier in the factorization.

The diagonal of \(R\) can provide information about numerical rank.

%%%%%%%%%%%%%%%%%%%%%%%%%%%%%%%%%%%%%%%%%%%%%%%%%%%%%%%%%%%%
\subsection{Eigenvalue Conditioning}
\label{subsec:eigenvalue-conditioning}
%%%%%%%%%%%%%%%%%%%%%%%%%%%%%%%%%%%%%%%%%%%%%%%%%%%%%%%%%%%%

Eigenvalues can themselves be sensitive to matrix perturbations.

For symmetric matrices, eigenvalues are relatively well behaved under
small perturbations.

For general nonsymmetric matrices, eigenvalue sensitivity can be much more
severe.

Thus eigenvalue algorithms must be interpreted together with problem
conditioning.

%%%%%%%%%%%%%%%%%%%%%%%%%%%%%%%%%%%%%%%%%%%%%%%%%%%%%%%%%%%%
\subsection{Symmetric Eigenvalue Problems}
\label{subsec:symmetric-eigenvalue-problems}
%%%%%%%%%%%%%%%%%%%%%%%%%%%%%%%%%%%%%%%%%%%%%%%%%%%%%%%%%%%%

If

\[
A=A^T,
\]

then all eigenvalues are real.

Furthermore, \(A\) has an orthonormal eigenvector basis:

\[
\boxed{
A
=
Q\Lambda Q^T.
}
\]

This structure permits especially stable and efficient algorithms.

%%%%%%%%%%%%%%%%%%%%%%%%%%%%%%%%%%%%%%%%%%%%%%%%%%%%%%%%%%%%
\subsection{Schur Decomposition}
\label{subsec:schur-decomposition}
%%%%%%%%%%%%%%%%%%%%%%%%%%%%%%%%%%%%%%%%%%%%%%%%%%%%%%%%%%%%

For a real or complex square matrix, a Schur decomposition gives

\[
\boxed{
A
=
QTQ^*,
}
\]

where \(Q\) is unitary and \(T\) is upper triangular in the complex case.

The diagonal entries of \(T\) are the eigenvalues of \(A\).

Schur form is central to practical eigenvalue computation.

%%%%%%%%%%%%%%%%%%%%%%%%%%%%%%%%%%%%%%%%%%%%%%%%%%%%%%%%%%%%
\subsection{Numerical Linear Algebra Workflow}
\label{subsec:nla-workflow}
%%%%%%%%%%%%%%%%%%%%%%%%%%%%%%%%%%%%%%%%%%%%%%%%%%%%%%%%%%%%

A practical workflow is:

\begin{enumerate}
    \item determine matrix dimensions;
    \item identify symmetry or positive definiteness;
    \item determine whether the matrix is dense or sparse;
    \item inspect scaling;
    \item estimate or diagnose conditioning;
    \item identify the required task;
    \item choose a factorization or iterative method;
    \item avoid unnecessary explicit inverses;
    \item use pivoting when appropriate;
    \item exploit symmetry and sparsity;
    \item monitor residuals;
    \item use backward-error reasoning;
    \item refine the solution if necessary;
    \item verify eigenpairs or least-squares residuals;
    \item compare numerical rank with application requirements;
    \item evaluate memory and computational cost.
\end{enumerate}

%%%%%%%%%%%%%%%%%%%%%%%%%%%%%%%%%%%%%%%%%%%%%%%%%%%%%%%%%%%%
\subsection{Choosing a Linear-System Solver}
\label{subsec:choosing-linear-system-solver}
%%%%%%%%%%%%%%%%%%%%%%%%%%%%%%%%%%%%%%%%%%%%%%%%%%%%%%%%%%%%

For a general dense square system, use:

\[
\boxed{
\text{LU with pivoting}.
}
\]

For symmetric positive-definite systems, use:

\[
\boxed{
\text{Cholesky}.
}
\]

For large sparse symmetric positive-definite systems, consider:

\[
\boxed{
\text{preconditioned CG}.
}
\]

For large sparse nonsymmetric systems, consider:

\[
\boxed{
\text{GMRES or another Krylov method}.
}
\]

%%%%%%%%%%%%%%%%%%%%%%%%%%%%%%%%%%%%%%%%%%%%%%%%%%%%%%%%%%%%
\subsection{Choosing a Least-Squares Solver}
\label{subsec:choosing-least-squares-solver}
%%%%%%%%%%%%%%%%%%%%%%%%%%%%%%%%%%%%%%%%%%%%%%%%%%%%%%%%%%%%

For a well-conditioned full-rank dense least-squares problem, QR is a
standard choice.

For rank-deficient or highly ill-conditioned problems, consider:

\[
\boxed{
\text{SVD}
}
\]

or

\[
\boxed{
\text{rank-revealing QR}.
}
\]

Normal equations can be faster in some settings but should be used with
awareness of their squared condition number.

%%%%%%%%%%%%%%%%%%%%%%%%%%%%%%%%%%%%%%%%%%%%%%%%%%%%%%%%%%%%
\subsection{Choosing an Eigenvalue Method}
\label{subsec:choosing-eigenvalue-method}
%%%%%%%%%%%%%%%%%%%%%%%%%%%%%%%%%%%%%%%%%%%%%%%%%%%%%%%%%%%%

For one dominant eigenpair, consider:

\[
\boxed{
\text{power iteration}.
}
\]

For an eigenvalue near a known shift, consider:

\[
\boxed{
\text{shifted inverse iteration}.
}
\]

For dense full-spectrum computations, practical methods are based on:

\[
\boxed{
\text{Hessenberg or tridiagonal reduction}
+
\text{QR-type iteration}.
}
\]

For large sparse problems, specialized Krylov eigensolvers are often
preferred.

%%%%%%%%%%%%%%%%%%%%%%%%%%%%%%%%%%%%%%%%%%%%%%%%%%%%%%%%%%%%
\subsection{Common Errors}
\label{subsec:nla-common-errors}
%%%%%%%%%%%%%%%%%%%%%%%%%%%%%%%%%%%%%%%%%%%%%%%%%%%%%%%%%%%%

\subsubsection{Computing \(A^{-1}\) to Solve \(A\mathbf{x}=\mathbf{b}\)}

Explicit inversion is usually unnecessary.

Use a factorization or iterative solve instead.

%%%%%%%%%%%%%%%%%%%%%%%%%%%%%%%%%%%%%%%%%%%%%%%%%%%%%%%%%%%%
\subsubsection{Ignoring Pivoting}

Gaussian elimination without pivoting can fail even for nonsingular
matrices or can produce severe numerical error.

Partial pivoting should normally be used for general dense systems.

%%%%%%%%%%%%%%%%%%%%%%%%%%%%%%%%%%%%%%%%%%%%%%%%%%%%%%%%%%%%
\subsubsection{Treating a Small Residual as Proof of a Small Error}

A small residual does not guarantee small forward error when

\[
\kappa(A)
\]

is large.

Conditioning matters.

%%%%%%%%%%%%%%%%%%%%%%%%%%%%%%%%%%%%%%%%%%%%%%%%%%%%%%%%%%%%
\subsubsection{Confusing Problem Conditioning and Solver Stability}

An ill-conditioned system remains sensitive even when solved by an
excellent algorithm.

A stable method cannot remove information lost in the problem itself.

%%%%%%%%%%%%%%%%%%%%%%%%%%%%%%%%%%%%%%%%%%%%%%%%%%%%%%%%%%%%
\subsubsection{Using Cholesky on an Indefinite Matrix}

Cholesky requires symmetric positive definiteness in its standard form.

Symmetry alone is not sufficient.

%%%%%%%%%%%%%%%%%%%%%%%%%%%%%%%%%%%%%%%%%%%%%%%%%%%%%%%%%%%%
\subsubsection{Forming \(A^TA\) Without Considering Conditioning}

The condition number satisfies

\[
\boxed{
\kappa_2(A^TA)
=
\kappa_2(A)^2
}
\]

for full-column-rank \(A\).

This can significantly reduce numerical accuracy.

%%%%%%%%%%%%%%%%%%%%%%%%%%%%%%%%%%%%%%%%%%%%%%%%%%%%%%%%%%%%
\subsubsection{Assuming Classical Gram--Schmidt Is Always Numerically Orthogonal}

Floating-point cancellation can cause loss of orthogonality.

Modified Gram--Schmidt or Householder QR is usually more reliable.

%%%%%%%%%%%%%%%%%%%%%%%%%%%%%%%%%%%%%%%%%%%%%%%%%%%%%%%%%%%%
\subsubsection{Using Determinants to Test Numerical Singularity}

Computing

\[
\det(A)
\]

is usually not a good numerical test for singularity or conditioning.

Factorizations and singular values are more informative.

%%%%%%%%%%%%%%%%%%%%%%%%%%%%%%%%%%%%%%%%%%%%%%%%%%%%%%%%%%%%
\subsubsection{Confusing Small Singular Values With Exact Zero}

A singular value may be mathematically nonzero but below a meaningful
numerical tolerance.

Numerical rank depends on scale and precision.

%%%%%%%%%%%%%%%%%%%%%%%%%%%%%%%%%%%%%%%%%%%%%%%%%%%%%%%%%%%%
\subsubsection{Applying CG to a General Nonsymmetric Matrix}

Standard conjugate gradient requires

\[
\boxed{
A=A^T
}
\]

and

\[
\boxed{
A\succ0.
}
\]

For general nonsymmetric systems, another Krylov solver is required.

%%%%%%%%%%%%%%%%%%%%%%%%%%%%%%%%%%%%%%%%%%%%%%%%%%%%%%%%%%%%
\subsubsection{Assuming Jacobi or Gauss--Seidel Always Converges}

Stationary iteration convergence depends on the iteration matrix.

A sufficient structural condition such as diagonal dominance should not be
confused with universal convergence.

%%%%%%%%%%%%%%%%%%%%%%%%%%%%%%%%%%%%%%%%%%%%%%%%%%%%%%%%%%%%
\subsubsection{Ignoring Preconditioning}

A theoretically convergent Krylov method may converge extremely slowly on
an ill-conditioned system.

A good preconditioner can be more important than changing the iterative
solver.

%%%%%%%%%%%%%%%%%%%%%%%%%%%%%%%%%%%%%%%%%%%%%%%%%%%%%%%%%%%%
\subsubsection{Ignoring Sparse Fill-In}

A sparse matrix does not guarantee a sparse factorization.

Poor ordering can produce large amounts of fill-in.

%%%%%%%%%%%%%%%%%%%%%%%%%%%%%%%%%%%%%%%%%%%%%%%%%%%%%%%%%%%%
\subsubsection{Using Eigenvalue Algorithms on the Characteristic Polynomial}

Forming and solving the characteristic polynomial numerically is usually
inferior to matrix-based eigenvalue methods.

%%%%%%%%%%%%%%%%%%%%%%%%%%%%%%%%%%%%%%%%%%%%%%%%%%%%%%%%%%%%
\subsubsection{Ignoring Scaling}

Large magnitude differences among rows, columns, or variables can create
avoidable numerical difficulties.

Scaling can improve practical solver behavior.

%%%%%%%%%%%%%%%%%%%%%%%%%%%%%%%%%%%%%%%%%%%%%%%%%%%%%%%%%%%%
\subsection{Applications}
\label{subsec:nla-applications}
%%%%%%%%%%%%%%%%%%%%%%%%%%%%%%%%%%%%%%%%%%%%%%%%%%%%%%%%%%%%

\begin{applicationbox}

\textbf{Optimization.}
Newton, interior-point, sequential quadratic programming, and least-squares
methods repeatedly solve linear systems.

\medskip

\textbf{Differential Equations.}
Finite-difference and finite-element discretizations produce large sparse
linear systems.

\medskip

\textbf{Statistics.}
Regression, covariance analysis, principal components, and multivariate
methods rely on QR, Cholesky, eigenvalue methods, and SVD.

\medskip

\textbf{Data Science.}
Low-rank approximations, dimensionality reduction, matrix completion, and
latent-factor models depend on matrix factorizations.

\medskip

\textbf{Engineering.}
Structural, circuit, control, and simulation models require repeated matrix
solves and eigenvalue calculations.

\medskip

\textbf{Networks.}
Graph Laplacians produce sparse linear algebra problems used in diffusion,
ranking, clustering, and network analysis.

\medskip

\textbf{Scientific Computing.}
Efficient matrix kernels form the computational foundation of many large
simulation codes.

\end{applicationbox}

%%%%%%%%%%%%%%%%%%%%%%%%%%%%%%%%%%%%%%%%%%%%%%%%%%%%%%%%%%%%
\subsection{Exercises}
\label{subsec:nla-exercises}
%%%%%%%%%%%%%%%%%%%%%%%%%%%%%%%%%%%%%%%%%%%%%%%%%%%%%%%%%%%%

\begin{exercise}[1]
Define a vector \(2\)-norm.
\end{exercise}

\begin{exercise}[1]
Define the matrix \(1\)-norm.
\end{exercise}

\begin{exercise}[1]
Define the matrix \(\infty\)-norm.
\end{exercise}

\begin{exercise}[1]
Define the Frobenius norm.
\end{exercise}

\begin{exercise}[1]
Define the condition number of a nonsingular matrix.
\end{exercise}

\begin{exercise}[1]
Define a linear-system residual.
\end{exercise}

\begin{exercise}[1]
Define LU factorization.
\end{exercise}

\begin{exercise}[1]
Define Cholesky factorization.
\end{exercise}

\begin{exercise}[1]
Define QR factorization.
\end{exercise}

\begin{exercise}[1]
Define the singular value decomposition.
\end{exercise}

\begin{exercise}[2]
Compute

\[
\|\mathbf{x}\|_1,
\qquad
\|\mathbf{x}\|_2,
\qquad
\|\mathbf{x}\|_\infty
\]

for a given vector.
\end{exercise}

\begin{exercise}[2]
Compute

\[
\|A\|_1
\]

and

\[
\|A\|_\infty
\]

for a given matrix.
\end{exercise}

\begin{exercise}[2]
Given \(A\) and an approximate solution \(\widehat{\mathbf{x}}\), compute
the residual.
\end{exercise}

\begin{exercise}[3]
Compare residual size with actual forward error in a small linear system.
\end{exercise}

\begin{exercise}[3]
Compute a condition number for a diagonal matrix.
\end{exercise}

\begin{exercise}[3]
Explain why a nearly singular matrix has a large condition number.
\end{exercise}

\begin{exercise}[2]
Perform Gaussian elimination on a \(3\times3\) system.
\end{exercise}

\begin{exercise}[2]
Perform back substitution.
\end{exercise}

\begin{exercise}[2]
Perform forward substitution.
\end{exercise}

\begin{exercise}[3]
Construct an LU factorization without pivoting.
\end{exercise}

\begin{exercise}[3]
Verify the factorization by multiplying \(LU\).
\end{exercise}

\begin{exercise}[3]
Perform one step of partial pivoting.
\end{exercise}

\begin{exercise}[3]
Construct

\[
PA=LU
\]

for a small matrix.
\end{exercise}

\begin{exercise}[3]
Explain why a very small pivot can cause numerical problems.
\end{exercise}

\begin{exercise}[2]
Determine whether a given symmetric matrix is positive definite.
\end{exercise}

\begin{exercise}[3]
Compute a \(2\times2\) Cholesky factorization.
\end{exercise}

\begin{exercise}[3]
Use Cholesky factors to solve a linear system.
\end{exercise}

\begin{exercise}[3]
Compare the leading arithmetic cost of LU and Cholesky factorization.
\end{exercise}

\begin{exercise}[2]
Verify that an orthogonal matrix preserves the Euclidean norm.
\end{exercise}

\begin{exercise}[3]
Apply Gram--Schmidt to two vectors.
\end{exercise}

\begin{exercise}[3]
Construct a QR factorization of a small matrix.
\end{exercise}

\begin{exercise}[3]
Verify that

\[
Q^TQ=I.
\]
\end{exercise}

\begin{exercise}[3]
Construct a Householder reflector that maps a vector to a coordinate axis.
\end{exercise}

\begin{exercise}[3]
Construct a Givens rotation that eliminates one vector component.
\end{exercise}

\begin{exercise}[3]
Solve a least-squares problem using QR.
\end{exercise}

\begin{exercise}[3]
Solve the same least-squares problem using normal equations.
\end{exercise}

\begin{exercise}[3]
Compare the condition numbers of \(A\) and \(A^TA\).
\end{exercise}

\begin{exercise}[2]
Compute the singular values of a diagonal matrix.
\end{exercise}

\begin{exercise}[3]
Construct the SVD of a simple \(2\times2\) matrix.
\end{exercise}

\begin{exercise}[3]
Determine the rank from its singular values.
\end{exercise}

\begin{exercise}[3]
Determine numerical rank using a specified tolerance.
\end{exercise}

\begin{exercise}[3]
Compute a pseudoinverse for a simple diagonal matrix.
\end{exercise}

\begin{exercise}[3]
Use the pseudoinverse to obtain a minimum-norm solution.
\end{exercise}

\begin{exercise}[3]
Construct a rank-one approximation using the leading singular triplet.
\end{exercise}

\begin{exercise}[3]
Compute the spectral-norm truncation error from the next singular value.
\end{exercise}

\begin{exercise}[2]
Perform one power-method iteration.
\end{exercise}

\begin{exercise}[3]
Apply several power iterations to a diagonalizable matrix.
\end{exercise}

\begin{exercise}[3]
Estimate the dominant eigenvalue using the Rayleigh quotient.
\end{exercise}

\begin{exercise}[3]
Explain how the eigenvalue ratio affects power-method convergence.
\end{exercise}

\begin{exercise}[3]
Perform one shifted inverse-iteration step.
\end{exercise}

\begin{exercise}[3]
Explain why a shift near an eigenvalue accelerates targeting.
\end{exercise}

\begin{exercise}[3]
Perform one unshifted QR iteration on a \(2\times2\) matrix.
\end{exercise}

\begin{exercise}[2]
Write the Jacobi iteration for a given linear system.
\end{exercise}

\begin{exercise}[2]
Write the Gauss--Seidel iteration for the same system.
\end{exercise}

\begin{exercise}[3]
Perform three Jacobi iterations.
\end{exercise}

\begin{exercise}[3]
Perform three Gauss--Seidel iterations.
\end{exercise}

\begin{exercise}[3]
Compare their observed convergence.
\end{exercise}

\begin{exercise}[3]
Construct the Jacobi iteration matrix.
\end{exercise}

\begin{exercise}[3]
Compute its spectral radius.
\end{exercise}

\begin{exercise}[3]
Determine whether the iteration converges for every initial guess.
\end{exercise}

\begin{exercise}[3]
Test a matrix for strict diagonal dominance.
\end{exercise}

\begin{exercise}[3]
Perform one SOR step for a specified \(\omega\).
\end{exercise}

\begin{exercise}[2]
Construct

\[
\mathcal{K}_3(A,\mathbf{r}_0).
\]
\end{exercise}

\begin{exercise}[3]
Perform one conjugate-gradient iteration.
\end{exercise}

\begin{exercise}[3]
Verify that two CG search directions are \(A\)-conjugate in exact
arithmetic for a small example.
\end{exercise}

\begin{exercise}[3]
Explain why CG requires symmetric positive definiteness.
\end{exercise}

\begin{exercise}[3]
Explain the role of the condition number in CG convergence.
\end{exercise}

\begin{exercise}[3]
Describe how GMRES minimizes residual norm over a Krylov subspace.
\end{exercise}

\begin{exercise}[3]
Construct a Jacobi preconditioner.
\end{exercise}

\begin{exercise}[3]
Explain the tradeoff in choosing a preconditioner.
\end{exercise}

\begin{exercise}[3]
Count the nonzeros in a sparse matrix and estimate sparse matrix-vector
multiplication cost.
\end{exercise}

\begin{exercise}[3]
Give an example of fill-in during sparse elimination.
\end{exercise}

\begin{exercise}[3]
Perform one iterative-refinement correction.
\end{exercise}

\begin{exercise}[3]
Explain the connection between backward stability and conditioning.
\end{exercise}

\begin{exercise}[4]
Derive standard perturbation bounds for linear systems.
\end{exercise}

\begin{exercise}[4]
Study growth factors in Gaussian elimination.
\end{exercise}

\begin{exercise}[4]
Study backward stability of LU factorization with partial pivoting.
\end{exercise}

\begin{exercise}[4]
Derive Cholesky factorization recursively.
\end{exercise}

\begin{exercise}[4]
Study stability of classical and modified Gram--Schmidt.
\end{exercise}

\begin{exercise}[4]
Derive Householder QR.
\end{exercise}

\begin{exercise}[4]
Study rank-revealing QR factorization.
\end{exercise}

\begin{exercise}[4]
Prove

\[
\kappa_2(A^TA)
=
\kappa_2(A)^2
\]

for full-column-rank \(A\).
\end{exercise}

\begin{exercise}[4]
Prove the Eckart--Young theorem.
\end{exercise}

\begin{exercise}[4]
Study perturbation theory for singular values.
\end{exercise}

\begin{exercise}[4]
Study eigenvalue conditioning for normal and nonnormal matrices.
\end{exercise}

\begin{exercise}[4]
Derive the power-method asymptotic convergence rate.
\end{exercise}

\begin{exercise}[4]
Study Rayleigh quotient iteration for symmetric matrices.
\end{exercise}

\begin{exercise}[4]
Study Hessenberg reduction using Householder reflectors.
\end{exercise}

\begin{exercise}[4]
Study shifted QR eigenvalue algorithms.
\end{exercise}

\begin{exercise}[4]
Prove the stationary-iteration condition

\[
\rho(G)<1.
\]
\end{exercise}

\begin{exercise}[4]
Study optimal SOR parameters for model problems.
\end{exercise}

\begin{exercise}[4]
Derive conjugate gradient from quadratic minimization.
\end{exercise}

\begin{exercise}[4]
Derive the classical CG error estimate.
\end{exercise}

\begin{exercise}[4]
Study Lanczos iteration.
\end{exercise}

\begin{exercise}[4]
Study Arnoldi iteration and GMRES.
\end{exercise}

\begin{exercise}[4]
Investigate incomplete LU preconditioning.
\end{exercise}

\begin{exercise}[4]
Study sparse ordering strategies and elimination graphs.
\end{exercise}

\begin{exercise}[4]
Investigate mixed-precision iterative refinement.
\end{exercise}

%%%%%%%%%%%%%%%%%%%%%%%%%%%%%%%%%%%%%%%%%%%%%%%%%%%%%%%%%%%%
\subsection{Conceptual Summary}
\label{subsec:nla-summary}
%%%%%%%%%%%%%%%%%%%%%%%%%%%%%%%%%%%%%%%%%%%%%%%%%%%%%%%%%%%%

Numerical linear algebra develops stable and efficient matrix algorithms.

A linear system is

\[
\boxed{
A\mathbf{x}
=
\mathbf{b}.
}
\]

Its conditioning is measured by

\[
\boxed{
\kappa(A)
=
\|A\|
\|A^{-1}\|.
}
\]

General dense systems are commonly solved using

\[
\boxed{
PA=LU.
}
\]

Symmetric positive-definite systems admit

\[
\boxed{
A=LL^T.
}
\]

Least-squares problems can be solved using

\[
\boxed{
A=QR.
}
\]

The singular value decomposition is

\[
\boxed{
A=U\Sigma V^T.
}
\]

It provides information about:

\[
\boxed{
\text{rank}
+
\text{conditioning}
+
\text{least squares}
+
\text{low-rank approximation}.
}
\]

Large sparse systems often use iterative methods such as

\[
\boxed{
\text{CG}
}
\]

and

\[
\boxed{
\text{GMRES},
}
\]

combined with

\[
\boxed{
\text{preconditioning}.
}
\]

Reliable numerical matrix computation therefore requires attention to

\[
\boxed{
\text{structure}
+
\text{conditioning}
+
\text{stability}
+
\text{sparsity}
+
\text{computational cost}.
}
\]

%%%%%%%%%%%%%%%%%%%%%%%%%%%%%%%%%%%%%%%%%%%%%%%%%%%%%%%%%%%%
\subsection{Section Connections}
\label{subsec:nla-connections}
%%%%%%%%%%%%%%%%%%%%%%%%%%%%%%%%%%%%%%%%%%%%%%%%%%%%%%%%%%%%

\chapterconnection{
Numerical Linear Algebra builds on the error, conditioning, and stability
ideas of Section~\ref{sec:numerical-analysis}. Matrix factorizations such
as LU, Cholesky, QR, and SVD provide the computational foundation for
least squares, eigenvalue problems, optimization, differential-equation
solvers, statistics, and data analysis. Iterative Krylov methods become
especially important for large sparse systems arising from scientific
models. The next section, Numerical Optimization, develops computational
algorithms for minimizing and maximizing nonlinear functions, including
line-search methods, Newton and quasi-Newton methods, constrained
optimization, and large-scale numerical optimization.
}

%%%%%%%%%%%%%%%%%%%%%%%%%%%%%%%%%%%%%%%%%%%%%%%%%%%%%%%%%%%%
% SECTION SOURCE TRACKING
%%%%%%%%%%%%%%%%%%%%%%%%%%%%%%%%%%%%%%%%%%%%%%%%%%%%%%%%%%%%

%------------------------------------------------------------
% 11.2 Numerical Linear Algebra
%------------------------------------------------------------
%
% Source basis:
%   - Standard numerical linear algebra.
%   - Standard direct and iterative matrix algorithms.
%   - Standard matrix-factorization and eigenvalue methods.
%
% Used for:
%   - dense and sparse matrices;
%   - vector and matrix norms;
%   - Frobenius norm;
%   - linear-system residuals;
%   - forward and backward error;
%   - matrix conditioning;
%   - Gaussian elimination;
%   - forward and back substitution;
%   - LU factorization;
%   - partial and complete pivoting;
%   - growth factors;
%   - arithmetic complexity;
%   - symmetric positive-definite matrices;
%   - Cholesky factorization;
%   - orthogonal matrices;
%   - QR factorization;
%   - classical and modified Gram--Schmidt;
%   - Householder reflections;
%   - Givens rotations;
%   - least-squares solution by QR;
%   - normal-equation conditioning;
%   - singular value decomposition;
%   - singular values;
%   - numerical rank;
%   - Moore--Penrose pseudoinverse;
%   - minimum-norm solutions;
%   - low-rank approximation;
%   - Eckart--Young theorem;
%   - eigenvalue computation;
%   - Rayleigh quotient;
%   - power method;
%   - inverse and shifted inverse iteration;
%   - Rayleigh quotient iteration;
%   - QR iteration;
%   - Hessenberg and tridiagonal reductions;
%   - direct and iterative solvers;
%   - matrix splitting;
%   - spectral-radius convergence;
%   - Jacobi;
%   - Gauss--Seidel;
%   - SOR;
%   - Krylov subspaces;
%   - conjugate gradient;
%   - GMRES;
%   - Arnoldi process;
%   - preconditioning;
%   - incomplete factorizations;
%   - sparse storage;
%   - fill-in and reordering;
%   - iterative refinement;
%   - mixed-precision refinement preview;
%   - scaling and equilibration;
%   - rank-revealing factorizations;
%   - Schur decomposition;
%   - applications and exercises.
%
% Notes:
%   - Numerical Optimization is developed next in Section 11.3.
%   - Numerical ODEs and PDEs follow in Sections 11.4 and 11.5.
%   - Scientific Computing in Section 11.7 develops broader computational
%     implementation concepts.
%
%%%%%%%%%%%%%%%%%%%%%%%%%%%%%%%%%%%%%%%%%%%%%%%%%%%%%%%%%%%%